\documentclass[11pt]{article}

\usepackage[margin=1.03in]{geometry}
\usepackage[T1]{fontenc}
\usepackage{lmodern}
\usepackage{microtype}
\usepackage{amsmath,amssymb,amsthm,mathtools}
\usepackage{booktabs}
\usepackage{enumitem}
\usepackage{xcolor}
\usepackage{listings}
\usepackage{xurl}
\usepackage[hidelinks]{hyperref}
\hypersetup{
  pdftitle={Strict negativity and non-vanishing of Chen--Larson hypergeometric coefficients: mathematics, formalization, and AI-assisted research},
  pdfauthor={Johannes Schmitt},
  pdfsubject={A detailed mathematical proof architecture, Lean 4 formalization, and documented AI-assisted research process for the Chen--Larson nonvanishing predictions (one corrected)},
  pdfkeywords={Lean 4, formal verification, strata of differentials,
               hypergeometric series, Ionel relation, AI-assisted mathematics}
}

\newtheorem{theorem}{Theorem}[section]
\newtheorem{proposition}[theorem]{Proposition}
\newtheorem{lemma}[theorem]{Lemma}
\newtheorem{corollary}[theorem]{Corollary}
\theoremstyle{definition}
\newtheorem{definition}[theorem]{Definition}
\theoremstyle{remark}
\newtheorem{remark}[theorem]{Remark}

\newcommand{\C}{\mathsf C}
\newcommand{\PE}{\mathsf E}
\newcommand{\Q}{\mathbb Q}
\newcommand{\e}{\mathrm e}
\newcommand{\floor}[1]{\left\lfloor #1\right\rfloor}
\newcommand{\ceil}[1]{\left\lceil #1\right\rceil}
\newcommand{\Lean}{\textsc{Lean}}
\newcommand{\code}[1]{\texttt{#1}}

\lstdefinestyle{leanstyle}{
  basicstyle=\ttfamily\small,
  columns=fullflexible,
  keepspaces=true,
  breaklines=true,
  frame=single,
  framerule=0.25pt,
  xleftmargin=0.5em,
  xrightmargin=0.5em,
  aboveskip=0.8em,
  belowskip=0.8em
}
\lstset{style=leanstyle}

% Human--AI interaction cards (cf. Feng et al.), dependency-free framed boxes.
% Usage: \begin{haicard}[<conversation url>]{<title>}{<model>} ... \end{haicard}
% With a url, the title becomes a coloured hyperlink to the conversation.
\definecolor{cardlink}{RGB}{0,70,148}
\newsavebox{\haibox}
\newenvironment{haicard}[3][]{%
  \par\medskip\noindent
  \begin{lrbox}{\haibox}%
  \begin{minipage}{\dimexpr\linewidth-2\fboxsep-2\fboxrule\relax}%
    \small
    {\sffamily\bfseries
      \def\haiurl{#1}%
      \ifx\haiurl\empty #2\else\href{#1}{\textcolor{cardlink}{#2}}\fi
      \hfill #3}\par\vspace{0.45em}%
}{%
  \end{minipage}%
  \end{lrbox}%
  \fbox{\usebox{\haibox}}%
  \par\medskip
}
% Field label inside a card.
\newcommand{\haifield}[1]{\textit{#1.}\enspace\ignorespaces}
% A short quoted prompt inside a card.
\newcommand{\haiprompt}[1]{\par\vspace{0.2em}\leftskip1.2em\rightskip1.2em\relax
  \small\itshape #1\par\leftskip0pt\rightskip0pt\vspace{0.2em}}

\title{Strict negativity and non-vanishing of Chen--Larson hypergeometric coefficients\\[0.45em]
\large A mathematical proof, Lean formalization, and documented AI-assisted research process}
\author{Johannes Schmitt}
\date{Revised version (25 June 2026)}

\begin{document}
\maketitle

\begin{abstract}
Chen and Larson study tautological classes on the strata of holomorphic
abelian differentials, where they predict a vanishing result.  Using known
tautological relations on the moduli space of curves, they reduce this
prediction to the non-vanishing of certain coefficients of a quotient of
hypergeometric generating series, treating the residue classes
$g\equiv0,2\pmod3$ and $g\equiv1\pmod3$ through two separate series; they verify
the non-vanishing by computer for small genus.

We prove it in general: in both cases the relevant coefficient is non-vanishing
for every genus and every stratum, and in the first case it has, more strongly,
a uniform strict sign.  Along the way we correct an error in the Chen--Larson
derivation of the $g\equiv1\pmod3$ series, where a constant of the underlying
relation of Ionel had inadvertently been changed.

The paper falls into two parts.  Part~I gives the mathematical proofs; Part~II
documents the Lean~4 and Mathlib formalization --- whose only non-standard trust
assumption is the compiler invoked by \code{native\_decide} for the large finite
computations --- together with the long-horizon, multi-model generative-AI
process that produced the proofs, exposed false intermediate routes, and
uncovered the error in the printed Proposition~5.2 relation noted above.
\end{abstract}

\tableofcontents

\part{Mathematical results and proofs}\label{part:mathematics}

\section{Introduction}

\subsection{Geometric context and main statements}

Chen and Larson study tautological classes on strata of holomorphic abelian
differentials.  In Section~5 of their work~\cite{ChenLarson}, they pull back a
relation among $\kappa$-classes on $\mathcal M_g$ to the projectivized stratum
$\mathcal P(\mu)$ associated with a positive partition
$\mu=(m_1,\ldots,m_n)$ of $2g-2$.  Here ``positive'' means that every
$m_i\ge1$, and $\eta$ denotes the tautological line-bundle class in their
convention.  The resulting scalar problems are expressed through
\begin{equation}\label{eq:CL-series-intro}
  \C(t)=\sum_{r\ge0}\frac{(6r)!}{(3r)!(2r)!}
                   \left(\frac{t}{72}\right)^r.
\end{equation}

For $g\equiv0,2\pmod3$, Chen--Larson Proposition~5.1 reduces the desired
vanishing of $\eta^a$, where $a=\floor{g/3}+1$, to the non-vanishing of
\[
  [t^a]\frac{\prod_{i=1}^n\C(t/(m_i+1))}{\C(t)^{2g-2+n}}.
\]
Our first theorem gives the stronger information that this coefficient always
has the same strict sign.

\begin{theorem}[Proposition~5.1 coefficient]\label{thm:main}
Let $g\ge2$ satisfy $g\equiv0$ or $2\pmod3$, and set
$a=\floor{g/3}+1$.  For every positive partition
$\mu=(m_1,\ldots,m_n)$ of $2g-2$,
\begin{equation}\label{eq:main-neg}
  [t^a]\frac{\prod_{i=1}^n\C(t/(m_i+1))}
                   {\C(t)^{2g-2+n}}<0.
\end{equation}
In particular, the coefficient in Chen--Larson Proposition~5.1 is nonzero.
\end{theorem}

The sign itself is not assigned a geometric meaning; it is a uniform
strengthening of the non-vanishing needed in the geometric reduction.

\begin{corollary}\label{cor:geometric}
Under the hypotheses of Chen--Larson Proposition~5.1, the class $\eta^a$
vanishes on $\mathcal P(\mu)$ for every positive partition $\mu$ when
$g\equiv0,2\pmod3$.
\end{corollary}

\begin{proof}
Apply Theorem~\ref{thm:main} to the scalar coefficient in the cited reduction.
\end{proof}

In the remaining class $g\equiv1\pmod3$, the same geometric vanishing rests on a
different Chen--Larson reduction --- their Proposition~5.2 --- built from a
derivative-corrected series rather than the quotient above; we first introduce
the notation it requires.

Write $g=3a-2$, so that
$a=\floor{g/3}+1\ge2$, and put
\[
  M=2g-2=6a-6,
  \qquad q_i=m_i+1,
  \qquad N=\sum_iq_i=M+n.
\]
Let $L(t)=\C'(t)/\C(t)$ and define
\begin{align}
 F_\mu(t)&=\frac{\prod_{i=1}^n\C(t/q_i)}{\C(t)^N},
 \label{eq:Fmu-intro}\\
 D_\mu^{\mathrm{cor}}(t)
 &=M-2\left(N-\sum_iq_i^{-1}\right)t
   -12t^2\left(NL(t)-\sum_iq_i^{-2}L(t/q_i)\right).
 \label{eq:Dcor-intro}
\end{align}
The leading constant $M=\kappa_0$ is the correction to the factor printed in
Chen--Larson Proposition~5.2, where it appears as~$1$.

\begin{theorem}[Corrected Proposition~5.2 coefficient]\label{thm:main52}
Let $g\ge4$ satisfy $g\equiv1\pmod3$, let $a=(g+2)/3$, and let
$\mu=(m_1,\ldots,m_n)$ be a positive partition of $2g-2$.  With
$F_\mu$ and $D_\mu^{\mathrm{cor}}$ defined by
\eqref{eq:Fmu-intro}--\eqref{eq:Dcor-intro},
\[
  [t^a]F_\mu(t)D_\mu^{\mathrm{cor}}(t)\ne0.
\]
For $a\ge14$, equivalently $g\ge40$, this coefficient is strictly negative.
\end{theorem}

Section~\ref{sec:prop52} derives the corrected factor from Ionel's relation,
proves that it agrees in degree~$a$ with a convenient marked expression, and
then proves the sign and the remaining finite non-vanishing cases.

\begin{corollary}\label{cor:geometric52}
Combining Theorems~\ref{thm:main} and~\ref{thm:main52} with the geometric
reductions of Chen--Larson Propositions~5.1 and~5.2, with the latter corrected
as in Section~\ref{sec:prop52}, gives
$\eta^{\floor{g/3}+1}=0$ on $\mathcal P(\mu)$ for every positive partition
$\mu$ and every genus $g\ge2$.
\end{corollary}

The algebraic-geometric reductions are imported from Chen and Larson, apart
from the corrected constant in Proposition~5.2.  The new work here is the
scalar power-series analysis and its formal verification.

\subsection{Proof architecture}\label{sec:outline}

After the geometric reduction, the problem is a uniform coefficient inequality
for a quotient of formal power series.  Two features are responsible for the
difficulty: the quotient depends on an arbitrary partition, and the sign must
hold uniformly as the genus grows.  The proof first removes the partition,
then separates a negative main term from positive corrections, and finally
uses different exact or analytic arguments in four numerical ranges.

\paragraph{Removing the partition.}
Write
\[
  \log\C(t)=\sum_{r\ge1}c_rt^r.
\]
Section~\ref{sec:coeff} proves $c_r>0$.  With $q_i=m_i+1\ge2$ and
$\sigma_r=\sum_iq_i^{-r}$, the numerator becomes
\[
  \prod_i\C(t/q_i)=\exp\!\left(\sum_{r\ge1}c_r\sigma_rt^r\right).
\]
Every coefficient of this exponential is nondecreasing in each $\sigma_r$.
Since $\sigma_r\le n2^{-r}\le(N/2)2^{-r}$, replacing each $\sigma_r$ by this
upper bound gives the coefficientwise majorant $\C(t/2)^{N/2}$.  After the
nonpositive convolution terms are discarded, the target coefficient is
bounded above by an explicit partition-free quantity $U_a(N)$.  It remains to
prove
\[
  U_a(N)<0\qquad(6a-7\le N\le12a-8).
\]

\paragraph{A negative term against a positive budget.}
Let $X_r(N)$ be the coefficient of $t^r$ in $\C(t)^{-N}$ normalized by
$Nc_r$, and let $Y_r(N)$ be the similarly normalized coefficient of
$\C(t/2)^{N/2}$.  Section~\ref{sec:majorant} rewrites the majorant as
\[
  \frac{U_a(N)}{Nc_a}=X_a(N)+\text{positive correction terms}.
\]
The proof therefore needs a negative upper bound for $X_a(N)$ and an upper
bound for the positive corrections.  The key negative estimate is the
\emph{sign lock}: whenever $m\ge361$ and $N\le(40/3)m$, the coefficient
$X_m(N)$ is negative with an explicit margin.  Applied with $m=k$, it removes
all convolution indices $k>0.9a$; applied with $m=a$, it supplies the main
negative term.

\paragraph{Keeping the two saddle estimates together.}
The surviving corrections are bounded without expanding their coefficients.
For
\[
  P_m(x)=\sum_{r=0}^m\frac{x^r}{r!},
\]
a finite saddle estimate bounds a coefficient of an exponential series by
$P_m$ at a rationally chosen radius.  The decisive product inequality is
\[
  P_m(x)P_n(y)\le P_{m+n}(x+y).
\]
It lets the denominator and numerator bounds be combined before their
normalization losses are incurred.  Section~\ref{sec:saddle} displays the full
calculation; this is exactly the point at which a plausible earlier route,
which estimated the two factors separately, lost a fatal dyadic gain.

\paragraph{Finite ranges and the infinite tail.}
For $a\le8$, every positive partition is enumerated and its exact rational
coefficient is checked.  For $9\le a\le60$, the majorant is checked exactly on
the integer rectangle, and for $61\le a\le400$ a proved dyadic-interval kernel
certifies the same inequalities.  For $a\ge401$, the sign lock and direct
saddle estimates reduce the positive side to a budget below $10^{-8}$.  The
strip through $a=2000$ is finite, the transition strip $2001\le a<3000$ is
handled by exact ratio certificates, and the range $a\ge3000$ follows from
uniform adjacent-term contractions and geometric reserves.

\paragraph{The corrected Proposition~5.2 coefficient.}
The correction changes only the constant term of the source factor, and in the
required degree it gives the identity
\[
  T_a^{\mathrm{cor}}=T_a^{\mathrm{old}}+(M-1)b_a.
\]
For $a\ge14$, the quotient coefficient $b_a$ is negative by the rectangular
form of Theorem~\ref{thm:main}.  The printed coefficient
$T_a^{\mathrm{old}}$ is checked by exact intervals for $14\le a\le149$ and is
proved negative for $a\ge150$ by a new Gamma-expectation argument.  That
argument represents a truncated coefficient sum by Gamma moments, uses
integration by parts and Jensen's inequality for a positive main margin, and
controls the Taylor remainder by an explicit four-term bound.  The corrected
coefficient is then negative.  The degrees $2\le a\le13$ are settled by exact
and modular finite certificates.

\subsection{Two contributions and two reading routes}

The manuscript is intentionally divided into two parts.  Part~I is the
mathematical paper: all notation, theorem statements, proof interfaces, finite
certificate claims, and analytic estimates needed for the conclusions are
stated there.  The finite certificates are presented as exact named lemmas,
with their proof methods described, while the full machine-checkable data live
in the associated Lean development.  No fact about the AI workflow is a
logical premise of the mathematics.

Part~II is a second, substantial account of formal verification and the
research process.  It
explains how the mathematical interfaces are represented in Lean, which finite
steps use native evaluation, exactly what must be trusted, and how a
human-orchestrated collection of generative-AI systems developed and audited
the argument.  Occasional pointers from Part~I to the correspondence appendices
are navigational only.  Conversely, Part~II presupposes the mathematical
architecture of Part~I rather than restating it.

The formalized scope is the complete scalar power-series content of
Theorems~\ref{thm:main} and~\ref{thm:main52}, over $\Q$, including the finite
certificates and the Gamma-tail argument.  The imported algebraic-geometric
reductions are not formalized.  Appendix~\ref{app:lean-statement} reproduces
the public theorem layer, and Section~\ref{sec:formalization} records the trust
boundary.

\section{The partition-free majorant}\label{sec:majorant}

This section makes the partition removal sketched in Section~\ref{sec:outline}
precise, ending at the majorant inequality $b_a(\mu)\le U_a(N)$.

Let $\mu=(m_1,\ldots,m_n)$ be a positive partition of $2g-2$ and put
\begin{equation}\label{eq:qN}
  q_i=m_i+1\ge2,
  \qquad
  N=\sum_{i=1}^nq_i=2g-2+n.
\end{equation}
Set
\begin{equation}\label{eq:bdef}
  b_a(\mu)=[t^a]\frac{\prod_i\C(t/q_i)}{\C(t)^N}.
\end{equation}
For $g\equiv0,2\pmod3$ and $a=\floor{g/3}+1$, one has
\[
  g=3a-3\quad\hbox{or}\quad g=3a-1.
\]
Since $1\le n\le2g-2$, every relevant value of $N$ lies in the enlarged
rectangle
\begin{equation}\label{eq:rectangle}
  6a-7\le N\le12a-8.
\end{equation}
The enlargement is convenient because it treats both residue classes at once.

Define
\begin{equation}\label{eq:AZQ}
\begin{split}
 A_r(\mu)&=[t^r]\prod_i\C(t/q_i),\\
 Z_r(N)&=[t^r]\C(t)^{-N},\\
 Q_r(N)&=[t^r]\C(t/2)^{N/2}.
\end{split}
\end{equation}
Here rational powers are interpreted through the exponential of the formal
logarithm.  Since the logarithmic coefficients $c_r$ are nonnegative,
coefficients of
\[
  \exp\left(\alpha\sum_{r\ge1}c_r2^{-r}t^r\right)
\]
are polynomials in $\alpha$ with nonnegative coefficients.  Writing the
numerator as $\prod_i\C(t/q_i)=\exp(\sum_rc_r\sigma_rt^r)$ with
$\sigma_r=\sum_iq_i^{-r}$, its coefficients $A_r(\mu)$ are therefore nondecreasing
in each $\sigma_r$.  Since $q_i\ge2$ and $n\le N/2$ give
$\sigma_r\le n2^{-r}\le\tfrac N2\,2^{-r}$, replacing each $\sigma_r$ by its upper bound
$\tfrac N2\,2^{-r}$ --- which turns the numerator into $\C(t/2)^{N/2}$ --- only
increases coefficients.  It follows that
\begin{equation}\label{eq:PleQ}
  0\le A_r(\mu)\le Q_r(N).
\end{equation}

Expanding \eqref{eq:bdef}, dropping nonpositive convolution terms, and using
\eqref{eq:PleQ} gives the partition-free majorant
\begin{equation}\label{eq:Udef}
 b_a(\mu)\le U_a(N):=Z_a(N)+Q_a(N)+
   \sum_{\substack{1\le k<a\\Z_k(N)>0}}Z_k(N)Q_{a-k}(N).
\end{equation}
Thus Theorem~\ref{thm:main} follows once we prove
\begin{equation}\label{eq:Uneg}
  U_a(N)<0
\end{equation}
throughout \eqref{eq:rectangle}, except for the finitely many values where we
instead verify $b_a(\mu)<0$ directly.

\subsection{Normalized form}

We pass to rescaled coefficients whose leading terms are the clean values
$X_1=-1$ and $Y_1=1$.  In this normalization the majorant splits transparently
into one negative term and two positive contributions --- the split the rest of
the proof exploits.  For $r\ge1$ define
\begin{equation}\label{eq:XY}
  X_r(N)=\frac{Z_r(N)}{Nc_r},
  \qquad
  Y_r(N)=\frac{2^rQ_r(N)}{(N/2)c_r},
\end{equation}
and set
\begin{equation}\label{eq:R}
  R_{k,a}=\frac{c_kc_{a-k}}{c_a}.
\end{equation}
For a series $E(t)=\exp(\sum_{r\ge1}\ell_rt^r)=\sum_{r\ge0}e_rt^r$, logarithmic
differentiation gives the recurrence $r\,e_r=\sum_{k=1}^rk\,\ell_ke_{r-k}$.
Applied to $\C(t)^{-N}$ and $\C(t/2)^{N/2}$ --- whose logarithmic coefficients
are $-Nc_r$ and $\tfrac N2c_r2^{-r}$ --- and normalized by~\eqref{eq:XY}, it
gives
\begin{align}
  X_1&=-1,&
  X_r&=-1-\frac Nr\sum_{k=1}^{r-1}kR_{k,r}X_{r-k},
  \label{eq:Xrec}\\
  Y_1&=1,&
  Y_r&=1+\frac N{2r}\sum_{k=1}^{r-1}kR_{k,r}Y_{r-k}.
  \label{eq:Yrec}
\end{align}
A direct calculation transforms \eqref{eq:Udef} into
\begin{equation}\label{eq:Unorm}
 \frac{U_a(N)}{Nc_a}
  =X_a(N)+2^{-a-1}Y_a(N)
   +\frac N2\sum_{\substack{1\le k<a\\X_k(N)>0}}
      R_{k,a}2^{-(a-k)}X_k(N)Y_{a-k}(N).
\end{equation}
The proof therefore consists of a uniform positive lower bound for $-X_a$
and an upper budget for the remaining two positive contributions.

\section{Logarithmic coefficients}\label{sec:coeff}

Everything downstream rests on the logarithmic coefficients $c_r$ of $\C$,
defined by $\log\C(t)=\sum_{r\ge1}c_rt^r$: their nonnegativity is what removed
the partition in Section~\ref{sec:majorant}, and the rational bounds derived
here drive every later estimate.  We establish both in turn.

The series $\C$ satisfies the second-order linear differential equation
\[
  6t^2\C''+(12t-1)\C'+\tfrac56\,\C=0,\qquad \C(0)=1,
\]
from which the bridge identifies $\C$ with $\exp(\sum_rc_rt^r)$ --- the two
series solve it with the same constant term.  Passing to the logarithmic
derivative $L=\C'/\C$ turns it into a first-order (Riccati) equation, quadratic
in $L$; equating coefficients there yields
\begin{equation}\label{eq:crec}
  c_1=\frac56,
  \qquad
  c_r=6(r-1)c_{r-1}
       +\frac6r\sum_{i=1}^{r-2}i(r-1-i)c_ic_{r-1-i}
       \quad(r\ge2).
\end{equation}
In particular, $c_r>0$.  Write
\begin{equation}\label{eq:ddef}
  c_r=6^r(r-1)!d_r.
\end{equation}
Then
\begin{equation}\label{eq:drec}
  d_r=d_{r-1}+\frac1r\sum_{i=1}^{r-2}
   \frac{d_id_{r-1-i}}{\binom{r-1}{i}}.
\end{equation}
The sequence $d_r$ is nondecreasing and begins with
$d_1=d_2=5/36$.

The formal proof uses rational bounds throughout.

\begin{lemma}[Uniform coefficient bounds]\label{lem:d-bounds}
For every $r\ge1$,
\begin{equation}\label{eq:d-bounds}
  \frac5{36}\le d_r\le\frac4{25}.
\end{equation}
Consequently,
\begin{equation}\label{eq:c-bounds}
 \frac5{36}6^r(r-1)!\le c_r\le
 \frac4{25}6^r(r-1)!.
\end{equation}
\end{lemma}

\begin{proof}
The lower bound follows immediately from \eqref{eq:drec}.  For the upper
bound, let $A_r=[t^r]\C(t)$ and
$F_r=A_r/(6^r(r-1)!)$.  One has $d_r\le F_r$ and
\[
  \frac{F_{r+1}}{F_r}=1+\frac{5}{36r(r+1)}.
\]
A rational telescoping argument bounds $F_r$ by
\[
  \frac{F_3}{1-(5/36)(1/3-1/r)}
  \le F_3\frac{108}{103}<\frac4{25}.
\]
No estimate involving $\pi$ is required.
\end{proof}

The reciprocal-binomial estimate
\begin{equation}\label{eq:binom-recip}
  \sum_{i=1}^{p-1}\binom pi^{-1}\le\frac4p
\end{equation}
combines with \eqref{eq:drec} and \eqref{eq:d-bounds} to give an increment
bound.

\begin{lemma}[Increment and ratio control]\label{lem:d-ratio}
For $r\ge2$,
\begin{equation}\label{eq:d-increment}
  d_r-d_{r-1}\le\frac{64/625}{r(r-1)}.
\end{equation}
If $1\le s<m$, then
\begin{equation}\label{eq:d-ratio}
  1-\frac{2304}{3125}\frac{s}{m(m-s)}
  \le\frac{d_{m-s}}{d_m}\le1.
\end{equation}
\end{lemma}

\begin{proof}
The nonlinear summand in \eqref{eq:drec} is at most
$(16/625)\binom{r-1}{i}^{-1}$.  Summing with
\eqref{eq:binom-recip} gives \eqref{eq:d-increment}.  Telescoping from
$m-s$ to $m$ yields
\[
 d_m-d_{m-s}\le\frac{64}{625}
      \left(\frac1{m-s}-\frac1m\right).
\]
Divide by $d_m\ge5/36$.  The upper ratio bound follows from monotonicity.
\end{proof}

A further consequence, used throughout the positive-part estimate, is
\begin{equation}\label{eq:R-bound}
 R_{k,a}\le
 \frac{576}{3125}
 \frac1{(a-1)\binom{a-2}{k-1}}
 \qquad(1\le k<a).
\end{equation}
Indeed, the $d$-factor contributes at most
$(4/25)^2/(5/36)=576/3125$, and the factorial quotient is exactly the
reciprocal denominator in \eqref{eq:R-bound}.

\section{Finite verification through \texorpdfstring{$a=400$}{a=400}}\label{sec:finite}

The majorant is not strong enough for the smallest values of $a$, so the
proof begins with direct finite verification.  The cut points $8$, $60$,
and $400$ below are computational thresholds --- each marks where one exact
method becomes infeasible and a cheaper one takes over --- not mathematically
distinguished values; the analytic argument begins at $a=401$, where its margin
first beats the positive budget (Section~\ref{sec:completion}).

\subsection{Positive partitions for \texorpdfstring{$a\le8$}{a<=8}}

For $g\le23$, equivalently $a\le8$ in the relevant residue classes, every
positive partition of $2g-2$ is generated and the exact rational coefficient
$b_a(\mu)$ is evaluated.  The largest enumeration is $p(44)=75175$
partitions at $g=23$.  The executable theorem checks strict negativity for
all of them.  Independent cardinality statements verify selected partition
counts, reducing the risk of an incomplete generator.

\subsection{Exact majorant certificate for \texorpdfstring{$9\le a\le60$}{9<=a<=60}}

For $9\le a\le60$, the recurrence tables for $c_r$, $Z_r(N)$, and $Q_r(N)$
are evaluated exactly over $\Q$.  This verifies \eqref{eq:Uneg} at all
$10764$ pairs in the rectangle.  As a regression anchor, the least-negative
corner is recorded exactly as
\begin{equation}\label{eq:corner}
 \frac{U_9(100)}{100c_9}
 =-\frac{13391635371339739}{213818836571652096}.
\end{equation}

\subsection{Verified dyadic intervals for \texorpdfstring{$61\le a\le400$}{61<=a<=400}}

The remaining $470220$ pairs are certified by a small executable interval
kernel.  A dyadic interval stores outward-rounded endpoints with a $192$-bit
mantissa.  The exact rational recurrences are mirrored term by term, and
separate soundness lemmas prove that the tables enclose the exact values.
For the conditionally included term in \eqref{eq:Udef}, the checker uses the
convex hull of the product interval and $\{0\}$; consequently it never needs
to decide the sign of $Z_k(N)$ from an interval approximation.

The finite check uses a dyadic interval kernel whose rational enclosure
semantics are proved independently.  Eight evaluation chunks cover the
$N$-range and together prove \eqref{eq:Uneg} for every $61\le a\le400$.
The $192$-bit precision matches the reference \emph{Arb}
(arbitrary-precision ball arithmetic) calculation used during discovery; the
formal trust boundary for this certification is described in
Section~\ref{sec:formalization}.


\begin{proposition}[Certified finite range]\label{prop:finite-range}
The following finite statements hold.
\begin{enumerate}[label=\textup{(\roman*)},leftmargin=2.2em]
\item For every relevant genus with $a\le8$ and every positive partition
$\mu$ of $2g-2$, one has $b_a(\mu)<0$.
\item For every pair of integers $a,N$ with $9\le a\le400$ and
$6a-7\le N\le12a-8$, one has $U_a(N)<0$.
\end{enumerate}
\end{proposition}

\begin{proof}
Part~\textup{(i)} is the exhaustive exact-rational partition computation
summarized in the first subsection.  Part~\textup{(ii)} is the union of the
exact recurrence calculation for $9\le a\le60$ and the sound dyadic-interval
calculation for $61\le a\le400$.  In the latter calculation, every arithmetic
operation is outward rounded and the enclosure invariant is proved by
induction along the recurrences, so a negative upper endpoint proves the exact
rational inequality.  The ranges are disjoint and exhaustive.
\end{proof}

\section{The effective sign lock}\label{sec:signlock}

For $a\ge401$, the leading coefficient $X_a(N)$ must be negative by more than
the entire positive budget.  The same estimate is also applied at the
convolution index $m=k$ to discard the upper tenth of the convolution range.
This section isolates the two mathematical interfaces on which that conclusion
rests: a normalized nonlinear residual envelope and a near/far audit for the
alternating expansion of $X_m(N)$.

\subsection{The nonlinear residual envelope}\label{subsec:envelope}

Write
\[
  H(t)=\sum_{r\ge2}c_rt^r,
  \qquad
  E^-_p(N)=[t^p]\exp\!\bigl(-NH(t)\bigr).
\]
The coefficient with one nonlinear block is $-Nc_p$.  The following proposition
quantifies the total contribution of every expansion term with at least two
blocks.

\begin{proposition}[Normalized nonlinear residual]\label{prop:residual-envelope}
Suppose that $m\ge361$, $N\ge1$, $N\le(40/3)m$, and $p\ge2m/3$.  Then
\begin{equation}\label{eq:residual-envelope}
  \left|\frac{E^-_p(N)}{-Nc_p}-1\right|\le\frac{66}{5m}.
\end{equation}
\end{proposition}

\begin{proof}
Expanding the exponential by the number of nonlinear blocks gives the exact
identity
\begin{equation}\label{eq:Eminus-block-expansion}
 E^-_p(N)=-Nc_p+
   \sum_{r=2}^{\floor{p/2}}\frac{(-N)^r}{r!}[t^p]H(t)^r;
\end{equation}
the upper limit occurs because every factor of $H$ has degree at least~$2$.
The coefficient bounds of Lemma~\ref{lem:d-bounds}, together with the
composition estimate for $[t^p]H^r$, bound the normalized absolute value of the
$r$th summand by
\begin{equation}\label{eq:Delta-term}
 \Delta_{p,N}^{(r)}=
 \frac{36}{5}\left(\frac4{25}\right)^r4^{r-1}N^{r-1}
 \frac{(p-2r+1)!}{r!(p-1)!}.
\end{equation}
Thus the left side of~\eqref{eq:residual-envelope} is at most
$\sum_{r=2}^{\floor{p/2}}\Delta_{p,N}^{(r)}$.  The first term and consecutive
quotient are
\begin{align}
 \Delta_{p,N}^{(2)}
   &=\frac{1152}{3125}\frac{N}{(p-1)(p-2)},
 \label{eq:Delta-first}\\
 \frac{\Delta_{p,N}^{(r+1)}}{\Delta_{p,N}^{(r)}}
   &=\frac{(16/25)N}
   {(r+1)(p-2r)(p-2r+1)}.
 \label{eq:Delta-ratio}
\end{align}
For $r\le p/4$, the quotient gives a geometric majorant beginning with
\eqref{eq:Delta-first}.  For $r>p/4$, the factorial quotient in
\eqref{eq:Delta-term} is bounded directly by
$(p-1)!/(p-2r+1)!\ge(2r-1)!$, after which the remaining terms again form a
geometrically decreasing tail.  Substituting $N\le(40/3)m$ and
$p\ge2m/3$ into these two bounds gives
\[
  \sum_{r=2}^{\floor{p/2}}\Delta_{p,N}^{(r)}
  \le\frac{66}{5m}.
\]
The denominator-cleared scalar inequalities used in this last summation are
recorded in Appendix~\ref{app:signlock-audit}; no asymptotic estimate is used.
\end{proof}

\subsection{The alternating decomposition}

Set
\[
  \zeta=\frac{5N}{36m}.
\]
The hypothesis $N\le(40/3)m$ gives $0\le\zeta\le50/27$.  Moreover,
$c_1=5/6$ and
\[
 \C(t)^{-N}=\exp(-Nc_1t)\exp(-NH(t)).
\]
Expanding the linear exponential gives
\begin{equation}\label{eq:signlock-exact}
  -X_m(N)=\sum_{s=0}^{m}\frac{(-Nc_1)^s}{s!}
           \left(-\frac{E^-_{m-s}(N)}{Nc_m}\right).
\end{equation}
For $s<m$, define
\[
  \Pi_s=\prod_{i=1}^{s}\frac1{1-i/m},
  \qquad D_s=\frac{d_{m-s}}{d_m},
  \qquad
  \varepsilon_p=\frac{E^-_p(N)}{-Nc_p}-1.
\]
The coefficient ratio $c_{m-s}/c_m$ then rewrites the $s$th summand as
\begin{equation}\label{eq:signlock-summand}
  \frac{(-\zeta)^s}{s!}\Pi_sD_s(1+\varepsilon_{m-s}).
\end{equation}
The factor $\Pi_s$ has first-order expansion
$1+e_1(s)/m$, where $e_1(s)=s(s+1)/2$.  The nonlinear residual has a matching
first-order endpoint contribution $-\zeta/m$.  We therefore define $w_s$ by
\begin{equation}\label{eq:ws-def}
 \Pi_sD_s(1+\varepsilon_{m-s})
  =1+\frac{e_1(s)-\zeta}{m}+w_s.
\end{equation}
This is the recentering that removes all first-order errors from the part of
the alternating sum that will be compared with an exponential.

Split~\eqref{eq:signlock-exact} at $s=\floor{m/3}$.  Proposition
\ref{prop:residual-envelope} applies throughout the near range because
$m-s\ge2m/3$.  We obtain
\begin{equation}\label{eq:signlock-three}
\begin{aligned}
 -X_m(N)={}&
 \underbrace{\sum_{0\le s\le\floor{m/3}}
   \frac{(-\zeta)^s}{s!}
   \left(1+\frac{e_1(s)-\zeta}{m}\right)}_{\mathcal B_m(\zeta)}\\
 &+\underbrace{\sum_{0\le s\le\floor{m/3}}
   \frac{(-\zeta)^s}{s!}w_s}_{\mathcal R_{\mathrm{near}}}
 +\underbrace{\sum_{\floor{m/3}<s\le m}
   \frac{(-Nc_1)^s}{s!}
   \left(-\frac{E^-_{m-s}(N)}{Nc_m}\right)}_{\mathcal R_{\mathrm{far}}}.
\end{aligned}
\end{equation}

\subsection{The near base and the error audit}

\begin{lemma}[Twelve-term near-base bound]\label{lem:near-base}
For $m\ge361$ and $0\le z\le50/27$,
\begin{equation}\label{eq:near-base-bound}
  \mathcal B_m(z)\ge\lambda\left(1-\frac2m\right),
  \qquad
  \lambda=\sum_{r=0}^{9}\frac{(-50/27)^r}{r!}
  =\frac{678107852315029}{4323713773987629}.
\end{equation}
\end{lemma}

\begin{proof}
Put
\[
 A_{12}(z)=\sum_{s=0}^{11}\frac{(-z)^s}{s!},
 \qquad
 C_{12}(z)=\sum_{s=0}^{11}\frac{(-z)^s}{s!}\bigl(e_1(s)-z\bigr).
\]
The twelve-term prefix is $A_{12}(z)+C_{12}(z)/m$.  After the affine change
$z=(50/27)u$, two exact Bernstein-basis expansions on $0\le u\le1$ have
nonnegative coefficients and prove
\[
 A_{12}(z)\ge\lambda,
 \qquad
 A_{12}(z)+\frac{C_{12}(z)}{361}
   \ge\lambda\left(1-\frac2{361}\right).
\]
The expression is affine in $1/m$, so the endpoint inequalities imply the same
bound for every $m\ge361$.  Starting with each even $s\ge12$, pair the terms
of indices $s$ and $s+1$.  Put $A_s=1+(e_1(s)-z)/m$.  Since
$A_{s+1}=A_s+(s+1)/m$ and $A_s\ge1$, the bracket remaining after factoring
out $z^s/s!$ satisfies
\[
 A_s-\frac{z}{s+1}A_{s+1}
 =A_s\left(1-\frac{z}{s+1}\right)-\frac zm
 \ge1-\frac{50}{27\cdot13}-\frac{50}{27\cdot361}>0.
\]
If the truncation leaves a final unpaired term, its index is even and the term
is nonnegative.  Hence the omitted alternating tail is nonnegative.
\end{proof}

\begin{proposition}[Near/far sign-lock audit]\label{prop:signlock-error-audit}
Suppose that $m\ge361$, $N\ge1$, and $N\le(40/3)m$, with
$\mathcal R_{\mathrm{near}}$ and $\mathcal R_{\mathrm{far}}$ as in
\eqref{eq:signlock-three},
\begin{equation}\label{eq:signlock-error-budget}
  |\mathcal R_{\mathrm{near}}|+|\mathcal R_{\mathrm{far}}|
  \le\frac{2215}{m^2}.
\end{equation}
\end{proposition}

\begin{proof}
Let $u_s=\Pi_s-1$ and $v_s=1-D_s$.  The product error is separated by the
exact inequality
\begin{equation}\label{eq:cross-product-bound}
 |(1+u_s)(1-v_s)(1+\varepsilon)-
   (1+u_s-v_s+\varepsilon)|
 \le u_s(v_s+|\varepsilon|)+v_s|\varepsilon|(1+u_s).
\end{equation}
The increment estimate~\eqref{eq:d-ratio} gives
$v_s\le(28/25)s/m^2$ on the near range.  The logarithmic product estimate for
$\Pi_s$ gives
\[
 0\le\Pi_s-1-\frac{e_1(s)}m,
 \qquad
 \log\Pi_s\le\frac{146}{125}\frac{e_1(s)}m
  <\frac{2237}{10000}s,
\]
that is, with decimal values $1.168$ and $0.2237$ respectively.  Finally, expansion~\eqref{eq:Eminus-block-expansion}
splits the recentered residual into the two endpoint terms of $H^2$, the
remaining two-block sum, and the terms with at least three blocks.  Applying
\eqref{eq:Delta-term} to those pieces and summing their Poisson weights gives
the following complete audit:
\begin{center}
\begin{tabular}{p{0.55\textwidth}r}
\toprule
source of the weighted error & bound\\
\midrule
$\Pi_s-1-e_1(s)/m$ & $426/m^2$\\
$d$-ratio drift $1-D_s$ & $13/m^2$\\
recentered endpoint part of the two-block term & $184/m^2$\\
non-endpoint two-block terms & $234/m^2$\\
terms with at least three nonlinear blocks & $573/m^2$\\
product cross terms from~\eqref{eq:cross-product-bound} & $784/m^2$\\
far signed tail & $1/m^2$\\
\midrule
sum & $2215/m^2$\\
\bottomrule
\end{tabular}
\end{center}
Each line is an exact rational inequality.  Appendix~\ref{app:signlock-audit}
records the block definitions, the geometric ratios, and the Bernstein or
endpoint certificate used for each line.  Summing the seven bounds proves
\eqref{eq:signlock-error-budget}.
\end{proof}

\begin{theorem}[Sign lock]\label{thm:signlock}
If $m\ge361$, $N\ge1$, and $N\le(40/3)m$, then
\begin{equation}\label{eq:signlock}
 X_m(N)\le-
 \left(\lambda\left(1-\frac2m\right)-\frac{2215}{m^2}\right)<0.
\end{equation}
\end{theorem}

\begin{proof}
Apply Lemma~\ref{lem:near-base} and
Proposition~\ref{prop:signlock-error-audit} to
\eqref{eq:signlock-three}.  The remaining scalar expression is positive at
$m=361$ by exact rational arithmetic, and
$m^2\lambda(1-2/m)$ is increasing for $m\ge361$.
\end{proof}

The value $9/10$ in the convolution cutoff is forced by the sign-lock
hypothesis rather than chosen for convenience.  Set
\[
  k_{\max}(a)=\floor{9a/10}.
\]
If $a\ge401$ and $k>k_{\max}(a)$, then $k\ge361$ and
\[
  N\le12a-8\le\frac{40}{3}k.
\]
Theorem~\ref{thm:signlock} gives $X_k(N)\le0$, so the positive sum in
\eqref{eq:Unorm} may be restricted to
\begin{equation}\label{eq:retained}
  1\le k\le\floor{9a/10}.
\end{equation}
For every retained index, $j=a-k\ge a/10$.

\section{Finite exponential prefixes and the direct saddle method}
\label{sec:saddle}

The positive corrections in~\eqref{eq:Unorm} involve a product of two
coefficients.  Estimating those coefficients separately and only then
multiplying loses precisely the dyadic gain needed later.  The purpose of this
section is therefore not merely to bound each factor, but to retain their
joint structure until the two finite exponential prefixes have been combined.

For $m\in\mathbb N$ and $x\ge0$, write
\begin{equation}\label{eq:exp-prefix}
  P_m(x)=\sum_{r=0}^m\frac{x^r}{r!}.
\end{equation}

\subsection{Generic finite saddle lemmas}

\begin{lemma}[Finite-prefix saddle bound]\label{lem:finite-saddle}
Let $m\in\mathbb N$, let $\rho\ge0$, and let
$\ell_1,\ldots,\ell_m\ge0$.  If
\[
 \mathcal E_m(\ell)=[t^m]\exp\left(\sum_{r=1}^{m}\ell_rt^r\right),
\]
then
\begin{equation}\label{eq:finite-saddle}
  \rho^m\mathcal E_m(\ell)
  \le P_m\left(\sum_{r=1}^m\ell_r\rho^r\right).
\end{equation}
\end{lemma}

\begin{proof}
After the change of variable $t\mapsto\rho t$, the left side is the
coefficient of $t^m$ in
\[
 \exp\left(\sum_{r=1}^m\ell_r\rho^rt^r\right).
\]
Expanding by the number $q$ of selected monomials, its $q$-block contribution
is at most
\[
 \frac{1}{q!}\left(\sum_{r=1}^m\ell_r\rho^r\right)^q.
\]
At most $m$ blocks can contribute to degree~$m$, because every block has
positive degree.  Summing
$q=0,\ldots,m$ gives~\eqref{eq:finite-saddle}.
\end{proof}

\begin{lemma}[Prefix convolution]\label{lem:prefix-conv}
For $m,n\in\mathbb N$ and $x,y\ge0$,
\begin{equation}\label{eq:prefix-conv}
  P_m(x)P_n(y)\le P_{m+n}(x+y).
\end{equation}
\end{lemma}

\begin{proof}
Expand the product over $0\le p\le m$, $0\le q\le n$, enlarge the rectangle
to the triangle $p+q\le m+n$, and group by $r=p+q$.  The binomial identity
then gives
\[
 \sum_{p+q=r}\frac{x^py^q}{p!q!}=\frac{(x+y)^r}{r!}.
\]
Every newly added summand is nonnegative.
\end{proof}

The same argument will be used in the two elementary forms
\begin{equation}\label{eq:prefix-absorptions}
  (1+d)^n\le P_n(nd),
  \qquad
  \frac{x^r}{r!}\le P_m(x)
  \quad(0\le r\le m;\,d,x\ge0).
\end{equation}

\subsection{Factorial-weighted exponent masses}

Since $c_r=6^r(r-1)!d_r$, the exponent masses in
Lemma~\ref{lem:finite-saddle} contain factorial weights.  We use the rational
saddle parameter
\begin{equation}\label{eq:beta}
  \beta=\frac{34}{15}.
\end{equation}
For $n\ge40$, put
\[
  u_{n,r}=\beta^r\frac{(r-1)!}{n^r}.
\]
Then
\begin{equation}\label{eq:u-ratio}
  \frac{u_{n,r+1}}{u_{n,r}}=\frac{34r}{15n}.
\end{equation}
The sequence is therefore first decreasing and then increasing, so every term
with $4\le r\le n$ is bounded by the sum of the two endpoints of that range.
The exact endpoint estimates are
\begin{equation}\label{eq:gas-components}
 n(u_{n,2}+u_{n,3})\le\frac16,
 \qquad n^2u_{n,4}\le\frac18,
 \qquad n^2u_{n,n}\le\frac12.
\end{equation}
For the last inequality, the quantity
$W_n=\beta^n n!/n^{n-1}$ satisfies $W_{40}\le1/2$ and
\[
  \frac{W_{n+1}}{W_n}
  =\frac{\beta}{(1+1/n)^{n-1}}\le1
  \qquad(n\ge40).
\]
Consequently,
\begin{equation}\label{eq:tempered-gas}
  \sum_{r=2}^n\beta^r\frac{(r-1)!}{n^r}\le\frac1n.
\end{equation}

A second estimate is sharper when the coefficient index is below a square-root
scale.  If $s\ge32$ and $2\le k\le s$, the consecutive quotient of
$(r-1)!/s^r$ shows that the middle terms are bounded by the low endpoint, and
one obtains
\begin{equation}\label{eq:small-gas}
  \sum_{r=2}^k\frac{(r-1)!}{s^r}
  \le\frac1{s^2}+\frac8{s^3}
  \le\frac5{4s^2}.
\end{equation}

\subsection{Single-factor saddle bounds}

Let $Z_k^+(N)$ be the coefficient of $t^k$ in
$\exp(N\sum_{r\ge1}c_rt^r)$.  Coefficientwise absolute values give
$Z_k(N)\le Z_k^+(N)$.  For $s=\ceil{\sqrt N}$,
Lemma~\ref{lem:finite-saddle}, Lemma~\ref{lem:d-bounds}, and
\eqref{eq:small-gas} give
\begin{equation}\label{eq:B-small}
  Z_k^+(N)\le(6s)^k
  P_k\left(\frac{5s}{36}+\frac15\right)
  \qquad(2\le k\le s).
\end{equation}

For indices at least~$40$, use the radii
\[
 \rho_Z=\frac{17}{45k}=\frac{\beta}{6k},
 \qquad
 \rho_Q=\frac{34}{45j}=\frac{\beta}{3j}.
\]
The tempered mass estimate~\eqref{eq:tempered-gas} yields
\begin{align}
 Z_k^+(N)&\le\left(\frac{45k}{17}\right)^k
 P_k\left(\frac{17N}{54k}+\frac{4N}{25k}\right),
 \label{eq:B-tempered}\\
 Q_j(N)&\le\left(\frac{45j}{34}\right)^j
 P_j\left(\frac{17N}{108j}+\frac{2N}{25j}\right).
 \label{eq:Q-tempered}
\end{align}

The normalization rests on
\begin{equation}\label{eq:factorial-lb}
  \left(\frac{25r}{68}\right)^r\le r!,
\end{equation}
which follows by induction from $(1+1/r)^r\le68/25$.  Together with the lower
bound in~\eqref{eq:c-bounds}, it implies
\begin{align}
 (6s)^k&\le\frac{36}{5}kc_k\frac{s^k}{k!},
 \label{eq:small-normalization}\\
 2^j\left(\frac{45j}{34}\right)^j
 &=\left(\frac{45j}{17}\right)^j
 \le\frac{36}{5}jc_j\left(\frac65\right)^j,
 \label{eq:Q-normalization}\\
 \left(\frac{45r}{17}\right)^r
 &\le\frac{36}{5}rc_r\left(\frac65\right)^r.
 \label{eq:tempered-normalization}
\end{align}
The identity
$(15/34)(68/25)=6/5$ explains the choice of~$\beta$.

\subsection{The direct product calculation}

\begin{theorem}[Direct product bounds]\label{thm:direct-product}
Let $a,k,j,N\in\mathbb N$ satisfy $k+j=a$, $1\le k<a$, and $N\le12a$.

If $s\ge32$, $2\le k\le s$, $N\le s^2$, and $j\ge40$, then
\begin{equation}\label{eq:small-core}
\begin{split}
 2\cdot 2^j Z_k(N)Q_j(N)
 \le{}&\frac{2592}{25}kjc_kc_j\\
 &\times P_{2a}\left(
 \frac{41}{36}s+\frac j5+
 \frac{29}{10}\frac aj+\frac15\right).
\end{split}
\end{equation}
If $k,j\ge40$, then
\begin{equation}\label{eq:tempered-core}
\begin{split}
 2\cdot 2^j Z_k(N)Q_j(N)
 \le{}&\frac{2592}{25}kjc_kc_j\\
 &\times P_{2a}\left(
 \frac a5+\frac{57}{10}\frac ak+
 \frac{29}{10}\frac aj\right).
\end{split}
\end{equation}
\end{theorem}

\begin{proof}
Because $Q_j(N)\ge0$, either $Z_k(N)\le0$, in which case both assertions are
immediate, or $0<Z_k(N)\le Z_k^+(N)$.  We treat the latter case.

For the small branch, abbreviate
\[
 A=\frac{5s}{36}+\frac15,
 \qquad
 B=\frac{17N}{108j}+\frac{2N}{25j}.
\]
Equations~\eqref{eq:B-small}, \eqref{eq:Q-tempered},
\eqref{eq:small-normalization}, and~\eqref{eq:Q-normalization} give
\begin{align*}
 2\cdot 2^j Z_k(N)Q_j(N)
 &\le 2\cdot 2^j Z_k^+(N)Q_j(N)\\
 &\le \frac{2592}{25}kjc_kc_j
      \frac{s^k}{k!}\left(\frac65\right)^jP_k(A)P_j(B).
\end{align*}
Now
\[
 \frac{s^k}{k!}\le P_k(s),
 \qquad
 \left(\frac65\right)^j
   =\left(1+\frac15\right)^j\le P_j(j/5).
\]
Applying Lemma~\ref{lem:prefix-conv} three times preserves the joint gain and
produces
\[
 P_k(s)P_k(A)P_j(j/5)P_j(B)
 \le P_{2a}\left(s+A+\frac j5+B\right).
\]
The rectangle estimate $N\le12a$ gives
\[
 B\le\left(\frac{17}{9}+\frac{24}{25}\right)\frac aj
   <\frac{29}{10}\frac aj,
\]
and $s+A=(41/36)s+1/5$.  This proves~\eqref{eq:small-core}.

For the tempered branch, set
\[
 A_k=\frac{17N}{54k}+\frac{4N}{25k},
 \qquad
 A_j=\frac{17N}{108j}+\frac{2N}{25j}.
\]
Using~\eqref{eq:B-tempered}, \eqref{eq:Q-tempered}, and
\eqref{eq:tempered-normalization} for both indices gives
\begin{align*}
 2\cdot 2^jZ_k(N)Q_j(N)
 &\le\frac{2592}{25}kjc_kc_j
       \left(\frac65\right)^{k+j}P_k(A_k)P_j(A_j)\\
 &=\frac{2592}{25}kjc_kc_j
       \left(\frac65\right)^aP_k(A_k)P_j(A_j).
\end{align*}
The residual power satisfies $(6/5)^a\le P_a(a/5)$, while prefix convolution
gives $P_k(A_k)P_j(A_j)\le P_a(A_k+A_j)$.  A final convolution therefore
bounds the product by $P_{2a}(a/5+A_k+A_j)$.  Finally,
\[
 A_k\le\frac{57}{10}\frac ak,
 \qquad
 A_j\le\frac{29}{10}\frac aj,
\]
again by $N\le12a$.  This proves~\eqref{eq:tempered-core}.  The factor
$2^j$ has remained inside~\eqref{eq:Q-normalization}; this is the dyadic gain
that disappears if the two normalized factors are bounded independently.
\end{proof}

\section{The positive budget for \texorpdfstring{$a\ge401$}{a>=401}}
\label{sec:positive}

We now insert Theorem~\ref{thm:direct-product} into the normalized majorant.
Let
\[
  N_-=6a-7,
  \qquad N_+=12a-8,
  \qquad s_-=\ceil{\sqrt{N_-}},
  \qquad s_+=\ceil{\sqrt{N_+}},
  \qquad j=a-k.
\]
For a rational $x<T$, define the certified exponential upper bound
\begin{equation}\label{eq:PE}
  \PE_T(x)=\sum_{r=0}^{T-1}\frac{x^r}{r!}
       +\frac{x^T}{T!}\frac1{1-x/T}.
\end{equation}

\begin{lemma}[Finite exponential majorant]\label{lem:finite-exp-majorant}
If $T\ge1$, $m\in\mathbb N$, and $0\le x<T$, then
\begin{equation}\label{eq:finite-exp-majorant}
  P_m(x)\le\PE_T(x).
\end{equation}
\end{lemma}

\begin{proof}
There is nothing to prove when $m<T$.  Otherwise, for $r\ge T$ the quotient
of consecutive summands $x^r/r!$ is $x/(r+1)\le x/T<1$.  Hence the part of
$P_m(x)$ beginning at degree~$T$ is bounded by the geometric series with first
term $x^T/T!$ and ratio $x/T$, which is exactly the final term in
\eqref{eq:PE}.
\end{proof}

For $401\le a\le2000$, define
\begin{align}
 \Xi_s(a,k)&=\frac{1139}{1000}s_++\frac j5+
             \frac{29}{10}\frac aj+1,
 \label{eq:Es}\\
 \Xi_t(a,k)&=\frac a5+\frac{57}{10}\frac ak+
             \frac{29}{10}\frac aj+2.
 \label{eq:Et}
\end{align}
For $401\le a\le2000$ and every retained $k$ one has $0\le \Xi_s(a,k)<800$.  In
the tempered regime $s_-<k$ --- and hence whenever $k>\ceil{\sqrt N}$ for an
admissible $N$, which is where $\Xi_t$ is used --- one also has
$0\le \Xi_t(a,k)<800$.  Thus each
occurrence of $\PE_{800}$ in \eqref{eq:PE} is well defined in the branch where
it is used.
The direct product theorem and elementary rational absorptions give
\begin{align}
 X_k(N)Y_j(N)
 &\le\frac{2581}{20}\frac{kj}{NN_+}\PE_{800}(\Xi_s(a,k)),
 &&k\le\ceil{\sqrt N},
 \label{eq:XY-small}\\
 X_k(N)Y_j(N)
 &\le\frac{2117}{20}\frac{kj}{N^2}\PE_{800}(\Xi_t(a,k)),
 &&k>\ceil{\sqrt N}.
 \label{eq:XY-tempered}
\end{align}
For the small branch, the comparison
\[
 \frac{2592}{25}N_+
 \le\frac53\frac{2581}{20}N
\]
uses the lower edge $N\ge N_-$; the factor $5/3$ is absorbed as
$P_1(2/3)$.  The numerical coefficient
$1139/1000$ exceeds $41/36$ by $1/9000$.  For the tempered branch,
$2592/25<2117/20$; the extra $+2$ in \eqref{eq:Et} is the explicit rational reserve used when the remaining prefix factors are absorbed.

Combining \eqref{eq:R-bound}, \eqref{eq:XY-small}, and
\eqref{eq:XY-tempered}, the retained positive convolution is bounded by a
sum of explicit terms
\begin{align}
 M_s(a,k)&=
 \frac{65}{N_+}\frac{kj}{(a-1)\binom{a-2}{k-1}}
 2^{-j}\PE_{800}(\Xi_s(a,k)),
 \label{eq:Ms}\\
 M_t(a,k)&=
 \frac{96}{N_-}\frac{kj}{(a-1)\binom{a-2}{k-1}}
 2^{-j}\PE_{800}(\Xi_t(a,k)).
 \label{eq:Mt}
\end{align}
The branch union must respect the regime conditions: the small bound
\eqref{eq:XY-small} applies only for $k\le\ceil{\sqrt N}$ and the tempered bound
\eqref{eq:XY-tempered} only for $k>\ceil{\sqrt N}$, and $\ceil{\sqrt N}$ moves
with $N$ across $s_-\le\ceil{\sqrt N}\le s_+$.  Outside its regime $M_t$ has no
content and can be large, so one may not simply take the unconditional maximum.
The theorem-facing edge term is therefore the regime-aware
\begin{equation}\label{eq:Medge}
  M_{\mathrm{edge}}(a,k)=
  \begin{cases}
    M_s(a,k), & k\le s_-,\\[1mm]
    \max\{M_s(a,k),\,M_t(a,k)\}, & s_-<k\le s_+,\\[1mm]
    M_t(a,k), & s_+<k,
  \end{cases}
\end{equation}
which dominates the retained summand for every admissible $N$: if
$k\le\ceil{\sqrt N}$ then $k\le s_+$ and the small bound, maximized at the upper
edge $N=N_+$, applies; if $k>\ceil{\sqrt N}$ then $s_-<k$ and the tempered bound,
whose prefactor is maximized at the lower edge $N=N_-$, applies; only in the
overlap $s_-<k\le s_+$ can the applicable branch depend on $N$, and there the
maximum of the two is taken.

For later statements, denote the complete positive side of
\eqref{eq:Unorm} by
\begin{equation}\label{eq:Vdef}
 \mathcal V_a(N)=2^{-a-1}Y_a(N)+\frac N2
   \sum_{\substack{1\le k<a\\X_k(N)>0}}
    R_{k,a}2^{-(a-k)}X_k(N)Y_{a-k}(N).
\end{equation}

\subsection{The finite strip \texorpdfstring{$401\le a\le2000$}{401<=a<=2000}}

\begin{proposition}[Finite positive budget]\label{prop:finite-positive-budget}
If $401\le a\le2000$ and $6a-7\le N\le12a-8$, then
\[
  \mathcal V_a(N)\le10^{-8}.
\]
\end{proposition}

\begin{proof}
For every retained index, the regime-aware term
$M_{\mathrm{edge}}(a,k)$ dominates the corresponding convolution summand.
An exact fixed-point calculation with outward rounding proves
\begin{equation}\label{eq:finite-edge-budget}
  \sum_{k\ \mathrm{retained}}M_{\mathrm{edge}}(a,k)
  \le\frac1{200000000}
\end{equation}
for all $401\le a\le2000$.  The rows are divided into sixteen blocks of one
hundred; each block evaluates the rational majorant~\eqref{eq:PE} at scale
$10^{20}$, and an induction over the arithmetic operations proves that the
stored integer is an upper bound for the exact rational value.

The solo contribution is bounded separately by the closed
factorial-composition estimate
\begin{equation}\label{eq:solo-budget}
  2^{-a-1}Y_a(N)\le\frac1{200000000}.
\end{equation}
Adding~\eqref{eq:finite-edge-budget} and~\eqref{eq:solo-budget} gives
\begin{equation}\label{eq:positive-target}
  \mathcal V_a(N)\le
  \frac1{200000000}+\frac1{200000000}=10^{-8}.
\end{equation}
\end{proof}

\subsection{The transition strip \texorpdfstring{$2001\le a<3000$}{2001<=a<3000}}

\begin{proposition}[Transition-strip positive budget]\label{prop:prefix-positive-budget}
If $2001\le a<3000$ and $6a-7\le N\le12a-8$, then
\[
  \mathcal V_a(N)\le10^{-8}.
\]
\end{proposition}

\begin{proof}
The direct product inequalities are applied pointwise with truncation cutoffs
$a$ and $8a$ in Lemma~\ref{lem:finite-exp-majorant}.  Exact adjacent-quotient
certificates cover the lower-tempered transition where a single symbolic
monotonicity statement would be inefficient.  Each band is bounded by its
boundary term times its proved geometric reserve, and the same
factorial-composition estimate as in~\eqref{eq:solo-budget} controls the solo
term.  The denominator-cleared inequalities and the finite list of boundary
rows are stated in Appendix~\ref{app:large-tail-details}; together they allocate
at most one half of $10^{-8}$ to the convolution and one half to the solo term.
\end{proof}

\subsection{The analytic tail \texorpdfstring{$a\ge3000$}{a>=3000}}

For the infinite range, the direct saddle bounds imply denominator-cleared
versions of~\eqref{eq:XY-small} and~\eqref{eq:XY-tempered}, with constants
$130$ and~$192$.  The reciprocal binomial factor is controlled by
\begin{equation}\label{eq:choose-power}
  \binom nk\ge\left(\frac nk\right)^k,
\end{equation}
and by a rational entropy-shadow inequality obtained from the maximal weighted
binomial term.  The retained range is split into a small branch and lower,
middle, and upper tempered bands.  Their candidate-majorant quotients satisfy
\begin{center}
\begin{tabular}{lccl}
\toprule
band & direction & quotient bound $\varrho$ & boundary term\\
\midrule
small ($k\le\ceil{\sqrt N}$) & forward & $\varrho\le\tfrac12$
  & first candidate term\\
tempered, lower/middle & forward & $\varrho\le(4a-1)/(4a)$
  & first term of the band\\
tempered, upper & reverse & $\varrho\le(4a-1)/(4a)$
  & $k=\floor{0.9a}$\\
\bottomrule
\end{tabular}
\end{center}
The tempered split occurs at $k=\floor{a/3}+10$.  Multiplication by the
corresponding reserve $1/(1-\varrho)$ keeps each band within its assigned
half-budget.  The actual guarded summand at $k=1$ is zero because $X_1=-1$;
the first small-branch term in the table is only the first term of the positive
candidate sequence.  The solo term is bounded through a $(10/7)^a$ envelope.

\begin{proposition}[Infinite positive budget]\label{prop:tail-positive-budget}
If $a\ge3000$ and $6a-7\le N\le12a-8$, then
\[
  \mathcal V_a(N)\le10^{-8}.
\]
\end{proposition}

\begin{proof}
Appendix~\ref{app:large-tail-details} proves the candidate domination, the
three adjacent-ratio estimates, the endpoint reserves, and the $(10/7)^a$ solo
bound as exact rational inequalities.  Proposition~\ref{prop:app-tail-final}
there assembles the four bands and the solo term into the displayed budget.
\end{proof}

\begin{theorem}[Positive envelope]\label{thm:positive-envelope}
For every integer $a\ge401$ and every integer $N$ with
$6a-7\le N\le12a-8$,
\begin{equation}\label{eq:envelope-result}
  \mathcal V_a(N)\le10^{-8}.
\end{equation}
\end{theorem}

\begin{proof}
Apply Proposition~\ref{prop:finite-positive-budget} when $a\le2000$,
Proposition~\ref{prop:prefix-positive-budget} when $2001\le a<3000$, and
Proposition~\ref{prop:tail-positive-budget} when $a\ge3000$.
\end{proof}

\section{Completion of the proof}\label{sec:completion}

\begin{proof}[Proof of Theorem~\ref{thm:main}]
For $a\ge401$, the rectangle gives $N\le12a-8<(40/3)a$, so
Theorem~\ref{thm:signlock}, applied with $m=a$, yields
\begin{equation}\label{eq:Xa-margin}
 X_a(N)\le-
 \left(\lambda\left(1-\frac2a\right)-\frac{2215}{a^2}\right).
\end{equation}
An exact endpoint comparison at $a=401$, followed by monotonicity, proves
\begin{equation}\label{eq:margin-target}
  10^{-8}<\lambda\left(1-\frac2a\right)-\frac{2215}{a^2}.
\end{equation}
Combining \eqref{eq:Unorm}, Theorem~\ref{thm:positive-envelope}, and
\eqref{eq:margin-target} gives $U_a(N)<0$ for every $a\ge401$.

Proposition~\ref{prop:finite-range} covers every remaining value $a\le400$.
\end{proof}

\section{The corrected Proposition~5.2}\label{sec:prop52}

The residue class $g\equiv1\pmod3$ requires a second scalar coefficient.  We
first derive the corrected source factor, then replace it in the required
degree by a marked factor whose constant term displays the correction.  The
resulting identity separates the problem into the already controlled quotient
coefficient and a new sign theorem for the series printed by Chen and Larson.
The latter is where the Gamma argument enters.

Throughout this section,
\begin{equation}\label{eq:prop52-notation}
 \begin{gathered}
  g=3a-2,\qquad M=2g-2=6a-6,\qquad q_i=m_i+1,\\
  N=\sum_iq_i=M+n,\qquad s_r=\sum_iq_i^{-r},\qquad
  F_\mu(t)=\frac{\prod_i\C(t/q_i)}{\C(t)^N},\\
  b_r(\mu)=[t^r]F_\mu(t),\qquad L(t)=\frac{\C'(t)}{\C(t)}.
 \end{gathered}
\end{equation}

\subsection{The source correction}\label{subsec:correction}

Ionel's relation~\cite[Proposition~1.7, equation~(1.13)]{Ionel}, with parameter
$b=1$, contributes in the target degree the factor
\begin{equation}\label{eq:ionel-b1}
 \kappa_0-2\kappa_1t-12\sum_{j\ge1}jc_j\kappa_{j+1}t^{j+1}.
\end{equation}
On $\mathcal M_g$, one has $\kappa_0=2g-2=M$.  Equation~(5.4) in the public
arXiv~v1 of Chen--Larson Proposition~5.2 instead has constant term~$1$.
Dividing Ionel's relation by $M$ would divide every term; changing only the
constant term is therefore not a normalization.

Under the stratum pullback used by Chen and Larson,
$\kappa_j$ for $j\ge1$ contributes $(N-s_j)\eta^j$, whereas
$\kappa_0$ remains~$M$.  Hence~\eqref{eq:ionel-b1} becomes
$D_\mu^{\mathrm{cor}}(\eta t)$, where
\begin{equation}\label{eq:Dcor52}
 D_\mu^{\mathrm{cor}}(t)
 =M-2(N-s_1)t-12t^2\left(NL(t)-\sum_iq_i^{-2}L(t/q_i)\right).
\end{equation}
The exponential prefactor is $F_\mu(\eta t)$, so the corrected geometric
reduction is
\begin{equation}\label{eq:corrected52}
  \eta^a[t^a]F_\mu(t)D_\mu^{\mathrm{cor}}(t)=0.
\end{equation}
The factor printed by Chen and Larson is~\eqref{eq:Dcor52} with the leading
$M$ replaced by~$1$; this is the only change, and it changes the coefficient.

\subsection{The marked factor and the correction identity}\label{subsec:identity52}

Define
\begin{equation}\label{eq:Kdef52}
 \Phi(t)=2t+12t^2L(t),
 \qquad
 K_\mu(t)=\sum_i(q_i-1)\Phi(t/q_i).
\end{equation}
The marked factors $M-K_\mu(t)$ and $1-K_\mu(t)$ are not equal to the source
factors as formal power series.  They may be substituted only after the
following degree-sensitive comparison.

\begin{lemma}[Source--marked bridge]\label{lem:bridge52}
For a scalar $\varkappa$, let $D_{\mu,\varkappa}$ be the expression in
\eqref{eq:Dcor52} with the leading $M$ replaced by $\varkappa$.  Then
\begin{equation}\label{eq:source-marked-series}
 F_\mu\bigl(D_{\mu,\varkappa}-(\varkappa-K_\mu)\bigr)
 =-2MtF_\mu+12t^2F_\mu'.
\end{equation}
Consequently, for every $r\ge1$,
\begin{equation}\label{eq:source-marked-coeff}
 [t^r]F_\mu D_{\mu,\varkappa}
 -[t^r]F_\mu(\varkappa-K_\mu)
 =\bigl(12(r-1)-2M\bigr)b_{r-1}(\mu).
\end{equation}
\end{lemma}

\begin{proof}
The constants $\varkappa$ cancel.  In the $t^2$ terms, the identity
$q_i^{-2}+(q_i-1)q_i^{-2}=q_i^{-1}$ combines the two sums.  The coefficient
of $t$ uses
$\sum_i(q_i-1)q_i^{-1}=n-s_1$ and $N-n=M$.  Therefore
\[
 D_{\mu,\varkappa}-(\varkappa-K_\mu)
 =-2Mt+12t^2\left(\sum_iq_i^{-1}L(t/q_i)-NL(t)\right).
\]
The parenthesis is $F_\mu'/F_\mu$, which proves
\eqref{eq:source-marked-series}.  Taking the coefficient of $t^r$ and using
$[t^r](t^2F_\mu')=(r-1)b_{r-1}(\mu)$ proves
\eqref{eq:source-marked-coeff}.
\end{proof}

At the Proposition~5.2 degree $r=a$, the multiplier in
\eqref{eq:source-marked-coeff} vanishes because $M=6a-6$.  We may therefore
define
\begin{align}
 T_a^{\mathrm{old}}(\mu)
   &:=[t^a]F_\mu D_{\mu,1}
     =[t^a]F_\mu(1-K_\mu),
 \label{eq:Told-def}\\
 T_a^{\mathrm{cor}}(\mu)
   &:=[t^a]F_\mu D_{\mu,M}
     =[t^a]F_\mu(M-K_\mu).
 \label{eq:Tcor-def}
\end{align}
Subtracting the two marked forms gives the central identity
\begin{equation}\label{eq:central52}
 \boxed{\ T_a^{\mathrm{cor}}(\mu)
 =T_a^{\mathrm{old}}(\mu)+(M-1)b_a(\mu).\ }
\end{equation}

\subsection{The two large-range sign inputs}\label{subsec:large52}

For a positive partition of $M=6a-6$, the parameter $N=M+n$ satisfies
\begin{equation}\label{eq:Nrange52}
  6a-5\le N\le12a-12.
\end{equation}
This interval lies inside the rectangle used in the Proposition~5.1 proof.

\begin{proposition}[Rectangular quotient sign]\label{prop:rectangle-bneg}
If $a\ge9$, $6a-7\le N\le12a-8$, and $\mu$ is a positive list with
$q_i=m_i+1$ and $\sum_iq_i=N$, then $b_a(\mu)<0$.
\end{proposition}

\begin{proof}
The coefficientwise argument of Section~\ref{sec:majorant} gives
$b_a(\mu)\le U_a(N)$ without using the value of $\sum_i(q_i-1)$.  For
$9\le a\le400$, Proposition~\ref{prop:finite-range} gives $U_a(N)<0$; for
$a\ge401$, the proof in Section~\ref{sec:completion} gives the same inequality
throughout the rectangle.
\end{proof}

In particular,~\eqref{eq:Nrange52} gives
\begin{equation}\label{eq:bneg52}
  b_a(\mu)<0\qquad(a\ge14).
\end{equation}
It remains to prove that the printed coefficient is negative as well.

\begin{proposition}[Printed coefficient in the middle range]
\label{prop:printed-mid}
For every $14\le a\le149$ and every positive partition $\mu$ of $6a-6$,
\[
  T_a^{\mathrm{old}}(\mu)<0.
\]
\end{proposition}

\begin{proof}
The coefficients $c_r$, $b_r(\mu)$, and $[t^r]K_\mu$ satisfy finite rational
recurrences.  A directed dyadic interval evaluation of these recurrences gives
a negative upper endpoint for every partition parameter cell in the stated
range.  The interval operations are outward rounded, and their enclosure
property is proved inductively from the rational recurrences.  Thus the finite
calculation proves the exact rational inequalities rather than importing a
floating-point assumption.  The certificate decomposition is listed in
Appendix~\ref{app:certs}.
\end{proof}

\subsection{The Gamma argument for the printed coefficient}
\label{subsec:gamma52}

We now prove the genuinely infinite sign input.  Define
\begin{equation}\label{eq:hk52}
 h_r=c_r(N-s_r),
 \qquad
 \tau_r=\sum_i\frac{m_i}{q_i^r},
 \qquad
 k_1=2\tau_1,
 \qquad
 k_r=12(r-1)c_{r-1}\tau_r\quad(r\ge2).
\end{equation}
Then
\[
 F_\mu(t)=\exp\left(-\sum_{r\ge1}h_rt^r\right),
 \qquad
 K_\mu(t)=\sum_{r\ge1}k_rt^r.
\]
Put
\begin{equation}\label{eq:low-split52}
 p=\floor{a/2},\qquad r_0=a-p-1,
\end{equation}
and form the low polynomials and series
\begin{equation}\label{eq:low-series52}
\begin{aligned}
 \Lambda(t)&=\sum_{r=1}^{p}h_rt^r,
 &J(t)&=\sum_{r=1}^{p}k_rt^r,\\
 E(t)&=\exp(-\Lambda(t))=\sum_{s\ge0}e_st^s,
 &W(t)&=E(t)(1-J(t))=\sum_{s\ge0}\omega_st^s.
\end{aligned}
\end{equation}
For $0\le s<a$, let
\begin{equation}\label{eq:gamma-weight52}
  \gamma_s=\frac{(a-s-1)!}{6^s(a-1)!},
  \qquad
  S_a(\mu)=\sum_{s=0}^{r_0}\gamma_s\omega_s.
\end{equation}
The weights are Gamma moments: if $G$ has the Gamma distribution of shape
$a$ and rate~$1$, and $T=1/(6G)$, then $\mathbb E_a[T^s]=\gamma_s$.

\begin{lemma}[Exact normalized split]\label{lem:printed-exact-split}
Define
\begin{align}
 H_a&=\sum_{s=0}^{r_0}\frac{h_{a-s}\omega_s}{Nc_a},
 &K_a&=\frac1{Nc_a}\sum_{s=0}^{r_0}k_{a-s}e_s,
 &\Omega_a&=\frac{\omega_a}{Nc_a}.
 \label{eq:HKOmega52}
\end{align}
Then
\begin{equation}\label{eq:printed-exact-split}
  -\frac{T_a^{\mathrm{old}}(\mu)}{Nc_a}=H_a+K_a-\Omega_a.
\end{equation}
\end{lemma}

\begin{proof}
Split both exponent and marked factor at degree~$p$.  Since
$2(p+1)>a$, a degree-$a$ term contains at most one high coefficient, whether
from the exponent or from the marked factor.  Every term containing such a
high exponent coefficient has complementary low degree at most
$r_0=a-p-1$.  Expanding the product and collecting the high exponent
term, the high marked term, and the residual low coefficient gives
\[
 T_a^{\mathrm{old}}
 =\omega_a-\sum_{s=0}^{r_0}h_{a-s}\omega_s
              -\sum_{s=0}^{r_0}k_{a-s}e_s.
\]
Division by $Nc_a$ and multiplication by $-1$ yields
\eqref{eq:printed-exact-split}.
\end{proof}

The next lemma isolates the analytic main term from three coefficient errors.

\begin{lemma}[Normalized split errors]\label{lem:printed-split-errors}
For $a\ge150$ and every positive partition of $6a-6$,
\begin{equation}\label{eq:printed-errors52}
 \left|H_a-S_a(\mu)\right|
   \le\frac{86069}{3125a^2},
 \qquad
 |K_a|\le\frac1{a^2},
 \qquad
 |\Omega_a|\le\frac1{a^2}.
\end{equation}
\end{lemma}

\begin{proof}
For the first estimate, compare
$h_{a-s}/(Nc_a)$ with the Gamma weight~$\gamma_s$.  The factorial part is
exactly $\gamma_s$, while the quotient $d_{a-s}/d_a$ is controlled by
Lemma~\ref{lem:d-ratio}; the remaining factor $1-s_{a-s}/N$ is between~$0$
and~$1$.  Summing the resulting weighted error and the low-degree endpoint
correction gives the constant
$86069/3125=36(2304/3125)+1$.

For $K_a$ and $\Omega_a$, replace the coefficients of $E$ and $W$ by the
positive coefficientwise majorants obtained from
$\exp(\sum_{r\le p}h_rt^r)(1+\sum_{r\le p}k_rt^r)$.  The bounds
\eqref{eq:c-bounds}, the partition inequalities $0\le s_r\le n$ and
$0\le\tau_r\le N$, and the factorial-mass estimates of
Section~\ref{sec:saddle} reduce both normalized tails to geometric series with
sum at most $a^{-2}$.  Appendix~\ref{app:gamma-tail-audit} records the exact
factorial sums and denominator-cleared endpoint inequalities.
\end{proof}

\begin{lemma}[Gamma main margin]\label{lem:gamma-main-margin}
For $a\ge150$ and every positive partition of $6a-6$,
\begin{equation}\label{eq:gamma-expectation-margin}
  \mathbb E_a\!\left[W\!\left(\frac1{6G}\right)\right]
  \ge\frac9{40(a-2)}.
\end{equation}
\end{lemma}

\begin{proof}
Integration by parts against the Gamma$(a,1)$ density gives
\begin{equation}\label{eq:gamma-ibp52}
 \mathbb E_a[W(T)]
 =\mathbb E_a\!\left[
 e^{-\Lambda(T)}\bigl(MT+6T^2\Lambda'(T)-J(T)\bigr)\right],
 \qquad T=\frac1{6G}.
\end{equation}
Indeed, differentiating $e^{-\Lambda(1/(6y))}$ produces
$6T^2\Lambda'(T)e^{-\Lambda(T)}$, while the logarithmic derivative of the
Gamma density contributes $MT=6(a-1)T$.  The boundary terms vanish at both
$0$ and $+\infty$, and the displayed integrands are absolutely integrable, so
the integration-by-parts identity is justified.

The bracket in~\eqref{eq:gamma-ibp52} has nonnegative coefficients.  Its
linear coefficient is $M-k_1\ge0$; for $1\le r<p$, the coefficient of
$t^{r+1}$ is
\[
  6r h_r-k_{r+1}\ge0,
\]
and the terminal coefficient of $t^{p+1}$ is $6ph_p\ge0$.  For $r=1$,
the sharper identity
\[
  6h_1-k_2-5M
  =5\sum_i\frac{(q_i-1)(q_i-2)}{q_i^2}\ge0
\]
gives $6h_1-k_2\ge5M$.  Therefore
\begin{equation}\label{eq:gamma-bracket-lower}
  MT+6T^2\Lambda'(T)-J(T)\ge5MT^2
  \qquad(T\ge0).
\end{equation}
Let $G_b\sim\mathrm{Gamma}(b,1)$ and set $T_b=1/(6G_b)$.  Changing the
Gamma shape by two gives
\begin{align}
 5M\,\mathbb E_a[T_a^2e^{-\Lambda(T_a)}]
 &=\frac{5M}{36(a-1)(a-2)}
   \mathbb E_{a-2}[e^{-\Lambda(T_{a-2})}]\\
 &=\frac5{6(a-2)}\mathbb E_{a-2}[e^{-\Lambda(T_{a-2})}].
 \label{eq:gamma-shape-shift}
\end{align}
Jensen's inequality now yields
$\mathbb E_{a-2}[e^{-\Lambda(T_{a-2})}]
 \ge e^{-\mathbb E_{a-2}[\Lambda(T_{a-2})]}$.  The coefficient bounds and Gamma
moments give the explicit mean estimate
\begin{equation}\label{eq:gamma-exponent-bound}
 \mathbb E_{a-2}[\Lambda(T)]
 \le \frac54+\frac5{2(a-3)}
       +\frac{96a^2}{25(a-5)^3}
 <\frac{13}{10}.
\end{equation}
Finally, $e^{-13/10}>27/100$.  Combining
\eqref{eq:gamma-ibp52}--\eqref{eq:gamma-exponent-bound} gives
\[
 \mathbb E_a[W(T)]
 \ge\frac5{6(a-2)}\frac{27}{100}
 =\frac9{40(a-2)}.
\]
\end{proof}

The finite sum $S_a(\mu)$ stops at $r_0$, so one still has to control a Taylor
remainder without interchanging a divergent collection of higher negative
Gamma moments.

\begin{lemma}[Gamma--Taylor truncation]\label{lem:gamma-truncation}
For $a\ge150$, let $s_0=\floor{a/8}$.  Then
\begin{equation}\label{eq:truncation-lower52}
 S_a(\mu)\ge
 \mathbb E_a\!\left[W\!\left(\frac1{6G}\right)\right]-R_a,
\end{equation}
where
\begin{equation}\label{eq:truncation-residue52}
 R_a=
 \frac{920}{2^{r_0+1}}
 +2\left(\frac56\right)^a
 +9\left(\frac9{10}\right)^{a-s_0}
 +920a\left(\frac3{10}\right)^{s_0+1}
 <\frac1{a^2}.
\end{equation}
\end{lemma}

\begin{proof}
Split the Gamma integral into $G\ge a/2$ and $G<a/2$.  On the first event,
$T\le x_1=1/(3a)$.  A positive coefficient majorant for $W$ at
$x_2=2/(3a)=2x_1$ is at most~$920$, so the Taylor terms beyond $r_0$ contribute
at most $920/2^{r_0+1}$.

On the lower event, a Gamma Chernoff estimate gives
$\mathbb P_a(G<a/2)\le(5/6)^a$.  Since
$|W(T)|=e^{-\Lambda(T)}|1-J(T)|
 \le e^{-\Lambda(T)}(1+J(T))\le2$,
the pointwise majorant contributes the second term of~$R_a$.
For the remaining Taylor comparison, split the monomial degrees at
$s_0=\floor{a/8}$.  The low degrees are bounded by the Gamma-moment ratio and
sum to at most $9(9/10)^{a-s_0}$; the high degrees are compared with the
$x_2$ majorant and sum to at most
$920a(3/10)^{s_0+1}$.  These are exactly the four terms in
\eqref{eq:truncation-residue52}.  Direct denominator clearing at the initial
values, followed by the monotone ratio bounds of
Appendix~\ref{app:gamma-tail-audit}, proves $R_a<a^{-2}$ for every $a\ge150$.
\end{proof}

\begin{proposition}[Printed coefficient in the Gamma range]
\label{prop:printed-tail}
For every $a\ge150$ and every positive partition $\mu$ of $6a-6$,
\begin{equation}\label{eq:printed-normalized-margin}
 -\frac{T_a^{\mathrm{old}}(\mu)}{Nc_a}
 \ge\frac9{40(a-2)}-\frac{95444}{3125a^2}>0.
\end{equation}
In particular, $T_a^{\mathrm{old}}(\mu)<0$.
\end{proposition}

\begin{proof}
Lemmas~\ref{lem:gamma-main-margin} and~\ref{lem:gamma-truncation} give
\[
  S_a(\mu)\ge\frac9{40(a-2)}-\frac1{a^2}.
\]
Insert this into the exact split of Lemma~\ref{lem:printed-exact-split} and
use the three estimates in Lemma~\ref{lem:printed-split-errors}.  The total
$a^{-2}$ coefficient is
\[
  1+\frac{86069}{3125}+1+1=\frac{95444}{3125},
\]
which proves the first inequality in~\eqref{eq:printed-normalized-margin}.
After denominators are cleared, positivity of the right side is equivalent to
\[
  28125a^2-3817760a+7635520>0.
\]
At $a=150$ the margin is
$3389201/20812500000$; writing the quadratic in powers of $a-150$ gives
positive coefficients, so the margin increases thereafter.
\end{proof}

Propositions~\ref{prop:printed-mid} and~\ref{prop:printed-tail} together give
\begin{equation}\label{eq:Toldneg52}
  T_a^{\mathrm{old}}(\mu)<0\qquad(a\ge14).
\end{equation}
Combining~\eqref{eq:central52}, \eqref{eq:bneg52}, and
\eqref{eq:Toldneg52} proves
\begin{equation}\label{eq:Tcorneg52}
  T_a^{\mathrm{cor}}(\mu)<0\qquad(a\ge14),
\end{equation}
because $M-1>0$.

\subsection{The remaining finite degrees}\label{subsec:finite52}

\begin{proposition}[Corrected coefficient in low degree]
\label{prop:corrected-low}
For $2\le a\le13$ and every positive partition $\mu$ of $6a-6$,
\[
  T_a^{\mathrm{cor}}(\mu)\ne0.
\]
\end{proposition}

\begin{proof}
The coefficient is computed from the finite recurrences as
\[
 T_a^{\mathrm{cor}}
   =Mb_a-\sum_{r=1}^{a}k_rb_{a-r}.
\]
For $2\le a\le8$, exhaustive partition enumeration evaluates this rational
number exactly.  For $9\le a\le13$, the same recurrence is evaluated modulo
the prime $1000000007$.  The denominator factors are proved invertible modulo
that prime, and every resulting residue is nonzero.  A rational number whose
denominator is invertible and whose reduction is nonzero cannot vanish.  The
partition-count and modular-certificate inventory appears in
Appendix~\ref{app:certs}.
\end{proof}

\begin{proof}[Proof of Theorem~\ref{thm:main52}]
Lemma~\ref{lem:bridge52} identifies the source coefficient
$[t^a]F_\mu D_\mu^{\mathrm{cor}}$ with $T_a^{\mathrm{cor}}(\mu)$, because
$12(a-1)-2M=0$.  Proposition~\ref{prop:corrected-low} proves non-vanishing for
$2\le a\le13$, and~\eqref{eq:Tcorneg52} proves strict negativity for
$a\ge14$.
\end{proof}

\begin{remark}[The correction changes the answer]
For $g=4$, so $a=2$ and $M=6$, and for $\mu=(1^6)$, one has
\[
 b_2=-\frac{195}{8},\qquad
 T_2^{\mathrm{old}}=\frac{45}{8},\qquad
 T_2^{\mathrm{cor}}=-\frac{465}{4}.
\]
Thus the printed and corrected coefficients even have opposite signs.  The
identity~\eqref{eq:central52} reads
$-465/4-45/8=5(-195/8)$ in this example.
\end{remark}

\begin{remark}[Scope]
The theorem is proved for positive integral partitions $m_i\ge1$.  The
corrected pullback factor makes sense in a broader rational-part setting, but
no global sign statement is asserted there.
\end{remark}

\part{Formal verification and AI-assisted research process}\label{part:formal-ai}

\section{Formalization architecture and trust}\label{sec:formalization}

Part~II treats formal verification and the AI-assisted research process as a
second, substantial contribution.  It is logically downstream from Part~I: the
formal development verifies the mathematical statements already formulated
there, and the process record explains how those statements and proofs were
generated, challenged, and revised.  Neither the Lean implementation nor the
provenance narrative is needed to follow the mathematical proofs.

The formalization nevertheless contributes more than a terminal correctness
check.  Its proof interfaces forced the analytic argument into explicit
finite-range, sign-lock, saddle, and tail modules; failed interfaces exposed
false global strategies; and the source-level formalization of the corrected
Proposition~5.2 separated a true coefficient statement from the geometric
statement that the printed factor was intended to encode.  This section maps
the mathematical architecture to Lean and states the resulting trust boundary.
Section~\ref{sec:ai} then records the AI-assisted workflow that produced and
audited both the mathematics and the formal development.

\subsection{Public statement layer}

The repository now has a concise public facade.  Its mathematical content is:
\begin{enumerate}[leftmargin=2em]
  \item the definition of the hypergeometric power series $\C$;
  \item the coefficient identity for $\C$;
  \item the quotient-series identity
  \[
    \C(t)^N F_\mu(t)=\prod_i\C(t/(m_i+1));
  \]
  \item strict negativity of the coefficient in Theorem~\ref{thm:main};
  \item the non-vanishing corollary.
\end{enumerate}
The bridge from the recurrence-defined $c_r$ to the formal power series is
proved by showing that both $\C$ and
$\exp(\sum_{r\ge1}c_rt^r)$ solve the same formal differential equation with
constant term one.

For the corrected Proposition~5.2 the public facade is
\code{Prop52/Theorem.lean}.  It imports \code{Prop52/Source.lean}, where the
source coefficient \code{sourceCoeff}, the coefficient-level product pinning,
the source--marked bridge \code{sourceCorrectedCoeff\_eq}, and the $g=4$
executable checks live.
The facade itself exposes only two genus-indexed source theorems: one for
non-vanishing and one for strict negativity.  The exact Lean identifiers are
listed in Appendix~\ref{app:lean-statement}.  The marked theorem names
\code{correctedCoeff\_nonvanishing} and \code{correctedCoeff\_neg} are retained
behind the facade in \code{Prop52/Assembly.lean}; the bridge turns them into the
source-shaped statement.  Both public theorems carry the same axiom report as
the Proposition~5.1 facade.

\subsection{Proof layers}

The production proof has the following conceptual dependency chain:
\begin{center}
\begin{tabular}{p{0.28\textwidth}p{0.62\textwidth}}
\toprule
layer & role\\
\midrule
series bridge & identifies the recurrence-defined coefficients with the
  Chen--Larson generating series\\
majorant & replaces the partition-dependent numerator by $Q_r(N)$ and proves
  $b_a(\mu)\le U_a(N)$\\
finite certificates & handles $a\le400$ by enumeration, exact rationals, and
  verified dyadic intervals\\
sign lock & eliminates $k>\floor{0.9a}$ and supplies the negative margin\\
direct saddle & proves the finite-prefix product inequalities
  \eqref{eq:small-core} and \eqref{eq:tempered-core}\\
positive assembly & checks the bounded edge budget and closes the prefix and
  infinite tail by ratio/reserve estimates\\
completion & combines all ranges into the unconditional public theorem\\
\bottomrule
\end{tabular}
\end{center}

\subsection{Trust boundary}

Almost all of the argument is checked by Lean's small logical kernel, the same
trusted core that underlies every Mathlib proof.  The exception is the family of
large finite computations --- the enumerations and the exact and interval
certificates --- which are run by \code{native\_decide}: Lean compiles the
finite Boolean check to machine code, executes it, and accepts the result.
This is faster by orders of magnitude than kernel reduction and is what makes
the millions of certificate cases feasible, but it widens the trusted base for
those steps to the Lean compiler and its runtime, together with any relevant
native implementations they invoke.  Following the Lean reference
manual~\cite{LeanAxioms}, each such use is recorded by a named axiom, and for
the present development the two additional native-evaluation axioms are
\code{Lean.ofReduceBool} and \code{Lean.trustCompiler}, alongside the standard
\code{propext}, \code{Classical.choice}, and \code{Quot.sound}.  Nothing else ---
in particular no \code{sorry}, no project-specific axiom, and no floating-point
assumption --- enters the final theorems.

This is verifiable in one command.  The repository ships a short audit script
that prints, for the series bridge, the majorant inequality, the
Proposition~5.1 negativity theorem and non-vanishing corollary, and the corrected
Proposition~5.2 theorems, the full list of axioms each depends on; a reader can
reproduce it from a clean checkout with
\begin{lstlisting}
lake exe cache get   # download the prebuilt Mathlib
lake build           # check every proof, including the certificates
lake env lean scripts/PublicAxiomsReport.lean   # print the axiom lists
\end{lstlisting}
The completed proof and the compact public facade are published as a tagged
repository release.  As an additional audit, the external checker
\code{lean4checker} can replay the ordinary (non-native) proof objects and can
detect certain bugs involving import state or declaration processing.  It does
not, however, remove the compiler trust attached to the native-evaluation
certificates: proofs that depend on \code{Lean.trustCompiler} cannot at present
be validated by an external proof checker without first replacing those
computations by kernel-checkable certificates~\cite{LeanAxioms}.

\section{AI-assisted research process and provenance}\label{sec:ai}

The AI-assisted research process is a substantive subject of this manuscript,
not merely a disclosure attached to the proof.  The record studies a concrete
case in which general-purpose generative systems proposed the central
mathematical architecture, generated and repaired a large Lean development,
performed adversarial audits, and identified an error in the printed source
relation.  The mathematical conclusions remain those of Part~I and stand
independently of this provenance; the purpose here is instead to make the
research method inspectable and to distinguish precisely among mathematical
generation, formal certification, and human orchestration.

The project began after Hannah Larson's talk \emph{Tautological classes on the
strata of differentials} in the Algebraic Geometry and Moduli seminar at ETH
Zurich on 10~June~2026, which described the conjecture of Chen and Larson.  The
author selected the problem and then worked almost entirely through
general-purpose AI tools.  The human role was orchestration: designing and
relaying prompts, maintaining the computational and Lean environment, deciding
from high-level progress signals when an independent review should be
solicited, curating the provenance record, and taking responsibility for the
final manuscript.  The author did not originate or modify a mathematical
lemma, estimate, proof strategy, or Lean proof.  The systems generated the
mathematical content and most of the formal development, but no theorem-level
claim was accepted merely because a model asserted it: acceptance rested on a
Lean proof or, for expository reformulations of already formalized statements,
on an independently checked mathematical derivation.

We describe the workflow through chronological \emph{human--AI interaction
cards}, following the format proposed by Feng et al.~\cite{FengAletheia}.  Each
card records the task, model, essential prompt, reply summary, and resulting
human action; its title links to the public conversation when one exists.  The
prompts are reproduced verbatim or in close paraphrase because prompt design
and the routing of work among independent systems are the most transferable
parts of the method.  The replies are summarized, with the most relevant
exchanges archived alongside the source.

In the content-based terminology of Feng et al., the mathematical generation
is classified as \textbf{Level~A}: no human supplied or edited the mathematical
lemmas, estimates, proof strategies, or Lean constructions.  This label does
not mean that the project was unattended.  The author selected the target,
organized the agents, interpreted only high-level signs of progress or
stalling, and relayed independent diagnoses.  In particular, the interventions
in Cards~6 and~7 were mathematical diagnoses produced by separate AI models
and forwarded to the worker without human alteration.  The classification is
therefore about the source of the mathematical content, not about the absence
of human project management.  We make no process-based claim about the
mathematical significance of the theorem: that assessment belongs to peer
review, and correctness rests on the proof rather than on the provenance.

The workflow had two phases.  During mathematical development, one
ChatGPT~5.5 Pro conversation was the central strand; independent conversations
alternated between adversarial audits and strategic guidance, and their
findings were fed back into that strand over ten drafts.  During formalization,
the work moved to Lean under a persistent autonomous goal.  When that run
stalled, two fresh-model reviews diagnosed defects in the current architecture
and supplied pivots that were relayed to the worker.  The cards below make this
routing, including the points of failure and correction, explicit.

\begin{haicard}[https://chatgpt.com/share/6a3661c0-131c-83ed-961d-1f82a69bfa0a]{Card 1: Central investigation strand}{ChatGPT 5.5 Pro}
\haifield{Task} Attack the Chen--Larson coefficient problem from scratch and
develop a candidate proof.\par
\haifield{Prompt} \haiprompt{Can you have a look at Conjecture 1.4 of
\textup{arXiv:2603.23850} (and its explicit reformulation Proposition 5.1) and
make a serious, sustained, high-effort research attack at trying to crack it?
This is an open problem, but in similar conversations you already managed to
crack several such problems, and you might have a perspective on this power
series that is not available easily to the authors. Would be really curious if
you can make progress here!}
\haifield{Reply (summary)} Proposed the coefficientwise majorization that
removes the partition, the normalized coefficients $X_r,Y_r$, and the strategy
of a negative leading term against small positive corrections.\par
\haifield{Continuation} This single conversation became the central strand: in
later turns it produced the \LaTeX{} write-up and Python verification scripts,
and carried the mathematics through ten successive drafts, fed by the audits and
guidance of Cards~2 and~3.\par
\end{haicard}

\begin{haicard}[https://chatgpt.com/share/6a36838a-3fa4-83ed-a43c-a3a5a1f1cde8]{Card 2: Adversarial mathematical audit}{ChatGPT 5.5 Pro}
\haifield{Task} Surface every defect in a draft package (notes and code) before
trusting it.\par
\haifield{Prompt} \haiprompt{Consider the appended zip file. Please read the
contained mathematical notes in detail and inspect all relevant code files, then
proceed with a full mathematical audit, going paragraph by paragraph to identify
any mathematical errors, gaps in arguments, hand-waving, misapplications of known
results, or mis-cited theorems --- anything that could affect the validity of the
presented arguments. Then give me a full report.}
\haifield{Reply (summary)} Returned concrete defects --- a far-tail constant
valid only in a smaller range, a positive-part bound vacuous at $a=401$, a
single-edge scan that missed the worst value of~$N$ --- each then repaired.\par
\haifield{Human action} Fed each reported defect back into the central strand for
repair, and re-ran the audit on the next draft.\par
\haifield{Technique} A fresh conversation used purely adversarially: an
exhaustive, paragraph-level audit of notes \emph{and} code, repeated each
revision.\par
\end{haicard}

\begin{haicard}[https://claude.ai/share/0ae838c7-7794-4868-9403-844c3e4db750]{Card 3: Strategic guidance}{Claude Fable 5}
\haifield{Task} Get higher-level direction, not line edits, when a draft stalled.\par
\haifield{Prompt} \haiprompt{I wanted to ask for advice about the appended pure
mathematics research note. For the theorem it tries to prove it shows some partial
progress, but gaps remain. Please look through it carefully, think hard, do your
own explorations or search for sources, and then give me a report focused on the
most helpful advice possible for a second AI which will be tasked to continue
pushing for a solution --- concrete pointers to results, obstructions, a
higher-level pivot, anything that could help crack the problem.}
\haifield{Reply (summary)} Returned strategic guidance --- which
effective-constant program to pursue, the worst-case regimes to target, and
higher-level pivots --- phrased as instructions the authoring model could act
on.\par
\haifield{Human action} Forwarded the guidance into the central strand to shape
the next revision.\par
\haifield{Technique} A second model used not to find errors but to suggest
direction, with the prompt explicitly asking for advice aimed at the next AI in
the loop.\par
\end{haicard}

\begin{haicard}{Card 4: First formalization push}{Claude Fable 5 (Claude Code)}
\haifield{Task} Begin the Lean development: the statement layer, the series
bridge, the partition-free majorant, the finite certificates, and the
sign lock.\par
\haifield{Prompt (excerpt)} \haiprompt{Following your advice, the author put a
tenth revision package together\dots{} this could be the basis of a first push
towards a formalization by you. Take the time to review any changes, then if you
think the basis is solid enough please go ahead and start developing the Lean
code here.\dots{} The goal would be to isolate the desired end theorem in a file
with as few imports, notation etc.\ as possible, with a clear certificate that
the proof is \textup{sorry}- and axiom-free.}
\haifield{Reply (summary)} Produced the bridge identifying $\C$ with
$\exp(\sum c_r t^r)$, the majorant inequality, the exact and dyadic-interval
finite certificates, and the rational toolkit underlying the sign lock.\par
\haifield{Human action} Reviewed and integrated the layers, then handed the
remaining analytic tail to a dedicated completion run.\par
\end{haicard}

\begin{haicard}{Card 5: Long-horizon completion drive}{OpenAI Codex CLI}
\haifield{Task} Drive the formalization to completion with minimal supervision.\par
\haifield{Prompt} \haiprompt{/goal Please continue working on the Lean
formalization of the argument, until it is either finished or until you hit a
serious obstacle which requires either mathematical input or technical
assistance beyond the standard difficulties and problem solving of the
formalization process.}
\haifield{Reply (summary)} Over many cycles of inspection, coding, compilation,
benchmarking, and repair, the worker proposed proof decompositions, searched
Mathlib, refactored, generated finite-certificate checkers, and profiled slow
computations.\par
\haifield{Human action} Monitored progress, spot-checked the Lean state, and
intervened only at the obstacle the goal itself anticipated (Card~6).\par
\haifield{Technique} A persistent goal that doubles as a completion
\emph{condition} attached to the thread~\cite{OpenAIGoals}; before this,
continuation had to be re-prompted every twenty to thirty minutes.\par
\end{haicard}

\begin{haicard}[https://chatgpt.com/share/6a36609a-3c80-83eb-91b8-209cd08f40ac]{Card 6: ``Is the worker stuck?'' --- a review of the process}{ChatGPT 5.5 Pro; Claude Opus 4.8}
\haifield{Task} Decide whether the multi-day Codex run was healthy exploration
or a dead end, and how to finish.\par
\haifield{Prompt (excerpt)} \haiprompt{Please review both the proof and the
formalization files in detail. (i)~Do you see any mathematical errors or
obstructions in the strategy, and are they fixable? (ii)~The formalization has
run with a Codex worker for about two days; is this the expected process of
investigation and optimization, or is the worker stuck in a local minimum or
dead end? (iii)~Any advice on how to finish, which could be forwarded to the
worker.}
\haifield{Reply (summary)} Run in parallel in both systems; the reviews found no
fatal obstruction in the main strategy but flagged unresolved analytic estimates
and a proliferation of proof interfaces, with the worker reshaping one
product-tail reduction rather than proving it, and advised freezing the
completion route around a single direct combined-product inequality.\par
\haifield{Human action} Forwarded the advice to the worker and gathered the
worker's own specific blocking questions for the deeper review of Card~7.\par
\haifield{Technique} Reviewing the \emph{process}, not only the mathematics, and
running the same review in two independent systems.\par
\end{haicard}

\begin{haicard}[https://chatgpt.com/share/6a36609a-3c80-83eb-91b8-209cd08f40ac]{Card 7: The correction that unstuck the proof}{ChatGPT 5.5 Pro; Claude Opus 4.8}
\haifield{Task} Resolve the worker's blocking questions about the large-tail
product estimate.\par
\haifield{Prompt (relayed, excerpt)} \haiprompt{[Relaying the worker's own
questions:] What is the most direct analytic proof of the sharp large-tail
product target, preferably by modifying the existing split rather than
introducing a large generated table?}
\haifield{Reply (summary)} The forwarded answer showed that the attempted route
--- bounding the two factors separately and multiplying --- is not merely
expensive but \emph{false} at the first large-tail cell $(a,k)=(3000,4)$, off by
about $3.26\cdot10^{557}$ because the split discards a factor $2^{a-k}$; it
prescribed the fix used in Section~\ref{sec:saddle}: keep the product together,
bound it with finite exponential prefixes, prove the prefix-convolution
inequality, and normalize only at the last step.\par
\haifield{Human action} Relayed the corrected route to the worker, which
implemented it as the direct-saddle method; the proof closed shortly
afterward.\par
\haifield{Technique} Escalating the worker's \emph{own} questions to a fresh,
stronger review, and letting a provably false target redirect the
architecture.\par
\end{haicard}

\paragraph{What the two interventions show.}
The episode in Cards~6 and~7 is worth keeping in the record.  The autonomous
worker had made genuine local progress around an interface that was globally
false: separating the two factors and multiplying them loses the dyadic
cancellation, and no amount of finite refinement could have repaired it.  What
exposed this was not a sharper estimate but the formal target itself, which
refused to close.  Formal verification thus did more than confirm a known proof:
it prevented a plausible but invalid analytic architecture from being mistaken
for a theorem, and it marked exactly the moment at which the autonomous worker
needed an outside perspective.  That perspective, too, came from an AI model; the
author's part was the orchestration --- reading a high-level signal that the run
had stalled and routing the question to a fresh review --- not any judgement
about the mathematics.  The mathematical content was AI-generated within a
human-orchestrated multi-agent process, and the human decisions that mattered
were about \emph{when} the worker was stuck, not how to fix it.

\paragraph{Scope of this record.}
The seven cards are a curated, not exhaustive, account: they record the
interactions that materially shaped the proof and omit many routine continuation
prompts and unproductive sessions.

The write-up itself was likewise AI-produced, in three stages.  The original
technical research note --- the tenth revision, focused purely on the
mathematics --- was authored by ChatGPT~5.5 Pro in the initial research strand of
Card~1.  After the formalization was complete, a fresh expository draft of the
present paper was produced by ChatGPT~5.5 Pro in a new conversation, following
the author's high-level instructions on organization and intended audience.  The
subsequent revisions --- including those prompted by an external review of the
release candidate --- were carried out in Claude Code under the author's
direction, together with further ChatGPT~5.5 Pro review.  As with the
mathematics, the author's contribution to the write-up was direction and
editorial judgement rather than composition; and acceptance of each formalized
claim rested on its Lean proof rather than on any model's assurance.

\subsection{Addendum: the corrected Proposition~5.2}\label{subsec:ai52}

The same methodology produced the companion result, with two features worth
recording: an AI audit caught a genuine error in the \emph{printed} (arXiv~v1)
relation, and the formalization then ran to completion without further
mathematical or technical intervention.

After the Proposition~5.1 result was shared with the authors --- who were open to
the approach and asked whether the remaining case $g\equiv1\pmod3$
(Proposition~5.2) could now be attacked --- the author posed that problem to a
ChatGPT~5.5 Pro
conversation,\footnote{\url{https://chatgpt.com/share/6a3bc8c9-e06c-83eb-b3e5-d098bb04573f}}
which produced a candidate non-vanishing proof for the \emph{printed}
Proposition~5.2 within three rounds.  A separate ChatGPT~5.5 Pro audit of that
candidate\footnote{\url{https://chatgpt.com/share/6a3bc91c-0460-83eb-857d-d821168e38b7}}
flagged that the derivation of the series in Proposition~5.2 was itself in error:
the constant $\kappa_0=2g-2$ of Ionel's relation had been replaced by $1$.  An
independent AI audit of the Chen--Larson paper
alone\footnote{\url{https://chatgpt.com/share/6a393ddd-bbc0-83eb-aa30-44806c9edc58}}
confirmed the error and supplied the corrected derivation of
Section~\ref{subsec:correction}; asked to validate the point, the authors
agreed, in private correspondence (June~2026), that the printed relation was in
error.  (As of June~2026, the public arXiv~v1 still prints the leading
constant~$1$.)

The corrected target was fed back into the original conversation, which --- in
two further rounds with intermediate
feedback\footnote{\url{https://chatgpt.com/share/6a3bc93b-3c68-83eb-8735-cb143fe10b12}}
--- assembled a complete bundle: an informal proof of the corrected statement, an
independent audit, the pinned Proposition~5.1 dependency, and a detailed Lean
formalization plan with milestones.  That bundle was handed to a fresh OpenAI
Codex CLI worker, which completed the entire Lean formalization --- the rectangle
export, the correction identity, the finite and modular certificates, the
in-Lean re-proof of the printed-coefficient sign, and the two public theorems
(sorry-free, with no project-specific axioms) --- \emph{autonomously, with no
further mathematical or technical
intervention}, in roughly twenty hours.  In the classification
of~\cite{FengAletheia} this is again Level~A, and the human role was smaller still
than for Proposition~5.1, reducing to relaying the problem, the audit verdict, and
the finished bundle.

Two points deserve emphasis.  First, the error was in the \emph{printed
manuscript}, not in any formalization: the series printed in Proposition~5.2 is a
valid power series, but it is not the one Ionel's relation produces, so a formal
proof about the printed coefficient would have certified a true statement that
does not carry the intended geometric meaning.  The defect was surfaced by
adversarial AI review of the paper itself --- exactly the kind of error formal
verification does not catch, because verification certifies \emph{what one
states}, while deciding what to state still needs independent mathematical
judgement.  Second, the contrast with Proposition~5.1 is instructive: given a
sufficiently detailed, AI-authored formalization plan, the autonomous worker
needed none of the mid-run redirection (Cards~6 and~7) that the first development
had required.

\paragraph{Quantitative notes.}
The Proposition~5.1 completion drive ran as a single Codex \texttt{/goal} thread,
for which the goal tracker recorded about sixty-four hours of elapsed time; the
separate Proposition~5.2 drive took roughly twenty hours by the same measure.
These are elapsed wall-clock figures as logged, not active compute time: they
include waiting and compilation and are not summed over parallel workers.  The
calendar span of the Proposition~5.1 thread was longer still --- it stayed open
from 15 to 20~June~2026 --- because the goal was paused between work sessions.
Human active time was not separately measured, and token counts and monetary
cost were not retained; the figures are given only to convey scale.  The formal
proof was checked under Lean~4.27.0 and a Mathlib revision pinned in the package
manifest.

\paragraph{On the persistent goal.}
The \texttt{/goal} prompt of Card~5 is reproduced as it was actually used.  By
later best-practice standards~\cite{OpenAIGoals} it is deliberately minimal: it
states the desired outcome and the condition under which to stop and ask for
help, but not the verification command, the permitted files, or an iteration
policy.  In practice the evolving Lean goal state, successful compilation, and
the axiom reports served as the implicit evidence surface; a reproduction attempt
should make those conditions explicit.

\paragraph{Limitations.}
Formal verification establishes the formal theorem relative to the trust boundary
of Section~\ref{sec:formalization}.  It does not by itself establish novelty or
significance, the correctness of bibliographic claims, or the faithfulness of
every sentence of the informal paper to the formal definitions; those remain
matters of human and community judgement.  The development also illustrates
characteristic failure modes of AI-assisted research, which this record
deliberately preserves: plausible but false intermediate estimates, a
proliferation of proof interfaces, and convincing local progress on a globally
false target.  Curated, redacted records of the command-line-agent sessions ---
every on-topic human instruction from the selected sessions, with the AI side
summarized --- together with the most relevant review exchanges are archived
under \code{docs/ai-correspondence/},
and the card titles link to the primary public conversations, so that the
provenance can be inspected directly.

\section{Reproducibility and repository organization}\label{sec:repro}

The repository is organized so that a reader can audit the public statements
without first traversing the development history.  Its four entry points are
\[
  \code{Prop51/Theorem.lean},\qquad \code{Prop52/Theorem.lean},\qquad
  \code{Prop51.lean},\qquad \code{Prop52.lean}.
\]
The two \code{Theorem.lean} files state the hypergeometric series, the
source-shaped quotient identities, and the public sign or non-vanishing
theorems.  The root files import the complete corresponding libraries.  The
paper-to-Lean map in Appendix~\ref{app:map} identifies the backend declaration
for every mathematical interface used in Part~I; Appendix~\ref{app:certs}
records the inputs and outputs of the finite certificates.

The archived release fixes the Lean toolchain at Lean~4.27.0 and pins the
Mathlib revision in the package manifest.  From a clean checkout the complete
verification and the public axiom audit are run by
\begin{lstlisting}
lake exe cache get
lake build
lake env lean scripts/PublicAxiomsReport.lean
\end{lstlisting}
The first command obtains the matching prebuilt Mathlib cache, the second
checks both proposition libraries and all finite certificates, and the third
prints the axiom dependencies of the public theorem layer.  The release also
contains the certificate-generation provenance and the curated AI
correspondence under \code{docs/ai-correspondence/}.  Thus the mathematical
proof, formal certificate boundary, and process record can be inspected
separately while remaining tied to the same versioned source tree.

\appendix

\section{Residual-envelope and sign-lock audit details}
\label{app:signlock-audit}

This appendix records the scalar estimates used in
Propositions~\ref{prop:residual-envelope} and
\ref{prop:signlock-error-audit}.  The purpose is not to reproduce the Lean
implementation, but to state the two analytic interfaces and the ordinary
inequalities used to certify them.  All constants below are rational, and every finite
polynomial assertion is certified after positive denominators have been
cleared.

\subsection{Summing the nonlinear residual blocks}

Recall the block majorant
\[
 \Delta_{p,N}^{(r)}=
 \frac{36}{5}\left(\frac4{25}\right)^r4^{r-1}N^{r-1}
 \frac{(p-2r+1)!}{r!(p-1)!}
 \qquad(2\le r\le p/2).
\]
Write
\[
 \Delta^{\mathrm{near}}_{p,N}
   =\sum_{2\le r\le p/4}\Delta_{p,N}^{(r)},
 \qquad
 \Delta^{\mathrm{far}}_{p,N}
   =\sum_{p/4<r\le p/2}\Delta_{p,N}^{(r)}.
\]
Under the sign-lock hypotheses, $m\ge361$, $2m\le3p$, and
$N\le(40/3)m$.  In particular,
\begin{equation}\label{eq:audit-pN}
 p\ge241,
 \qquad N\le20p.
\end{equation}
For $2\le r<p/4$, the exact quotient from
\eqref{eq:Delta-ratio} is bounded by
\[
 \frac{\Delta_{p,N}^{(r+1)}}{\Delta_{p,N}^{(r)}}
 \le \frac{64R}{75p},
 \qquad R=20.
\]
The right side is $256/(15p)<1$; in fact it is at most $1/11$ for
$p\ge241$.  Combining this with the first block gives
\begin{align}
 \Delta^{\mathrm{near}}_{p,N}
 &\le
 \frac{1152}{3125}\frac{20p}{(p-1)(p-2)}
 \frac1{1-256/(15p)}                                      \label{eq:audit-near-closed}\\
 &\le \frac{196416}{15625m}.                              \label{eq:audit-near-final}
\end{align}
For the last inequality, use $m\le3p/2$, first to obtain
\[
 20mp\le31(p-1)(p-2),
\]
and then use $(1-256/(15p))^{-1}\le11/10$.  These are
polynomial inequalities on $p\ge241$.

For the far blocks, the factorial quotient satisfies
\[
 \frac{(p-2r+1)!}{(p-1)!}\le\frac1{(2r-1)!}
 \qquad(p/4<r\le p/2).
\]
Together with $N\le20p<80r$, this gives
\begin{equation}\label{eq:audit-far-term}
 \Delta_{p,N}^{(r)}\le
 \widehat\Delta_r:=
 \frac95\left(\frac{16}{25}\right)^r
 \frac{(80r)^{r-1}}{r!(2r-1)!}.
\end{equation}
The consecutive quotient is
\begin{equation}\label{eq:audit-far-ratio}
 \frac{\widehat\Delta_{r+1}}{\widehat\Delta_r}
 =\frac{256}{5}\,
   \frac{(1+1/r)^{r-1}}{(2r)(2r+1)}
 \le\frac12
 \qquad(r\ge34),
\end{equation}
where $(1+1/r)^r<\mathrm e<68/25$ is sufficient for the last
inequality.  The first far index is
$r_* =\floor{p/4}+1\ge61$.  Applying the factorial lower bounds at this
index and then summing the geometric tail gives
\begin{equation}\label{eq:audit-far-final}
 \Delta^{\mathrm{far}}_{p,N}\le\frac1{1000m}.
\end{equation}
The endpoint inequality behind~\eqref{eq:audit-far-final} is stronger than
needed: after multiplication by $m$, its left side decreases with $p$ once
$p\ge241$, so it is enough to clear denominators at the first admissible
value.

Adding~\eqref{eq:audit-near-final} and
\eqref{eq:audit-far-final} yields
\[
 \sum_{r=2}^{\floor{p/2}}\Delta_{p,N}^{(r)}
 \le\left(\frac{196416}{15625}+\frac1{1000}\right)\frac1m
 <\frac{66}{5m},
\]
which is the last step in Proposition~\ref{prop:residual-envelope}.

\subsection{The seven sign-lock budgets}

We next make explicit how the error $w_s$ in~\eqref{eq:ws-def} is divided.
For $p=m-s$, write the normalized nonlinear expansion as
\[
 \varepsilon_p=\sum_{r=2}^{\floor{p/2}}\varepsilon_p^{(r)},
 \qquad
 \varepsilon_p^{(r)}
   =-\frac{(-N)^r}{r!Nc_p}[t^p]H(t)^r.
\]
The two-block term is separated into its endpoint compositions
$(2,p-2)$ and $(p-2,2)$ and the remaining compositions.  Thus
\[
 \varepsilon_p
 =\varepsilon_p^{(2,\mathrm{end})}
  +\varepsilon_p^{(2,\mathrm{mid})}
  +\varepsilon_p^{(\ge3)},
 \qquad
 \varepsilon_p^{(2,\mathrm{end})}
 =-N\frac{c_2c_{p-2}}{c_p}.
\]
The endpoint term is the source of the first-order correction
$-\zeta/m$ in~\eqref{eq:ws-def}.  Let
\[
 u_s=\Pi_s-1,
 \qquad v_s=1-D_s,
 \qquad z_*=\frac{50}{27}.
\]
After adding and subtracting $e_1(s)/m$ and $-\zeta/m$, the absolute value of
$w_s$ is bounded by six nonnegative contributions: the second-order residual
of~$\Pi_s$, the drift~$v_s$, the recentered endpoint two-block term, the
middle two-block term, the terms with at least three blocks, and the product
cross terms in~\eqref{eq:cross-product-bound}.  With the Poisson weight
$z_*^s/s!$, their complete sums satisfy
\begin{align}
 \sum_{s\le m/3}\frac{z_*^s}{s!}
   \left|\Pi_s-1-\frac{e_1(s)}m\right|
   &\le\frac{426}{m^2},                                      \label{eq:audit-P1}\\
 \sum_{s\le m/3}\frac{z_*^s}{s!}(1-D_s)
   &\le\frac{13}{m^2},                                       \label{eq:audit-P2}\\
 \sum_{s\le m/3}\frac{z_*^s}{s!}
   \left|\varepsilon_{m-s}^{(2,\mathrm{end})}+\frac\zeta m\right|
   &\le\frac{184}{m^2},                                      \label{eq:audit-P3a}\\
 \sum_{s\le m/3}\frac{z_*^s}{s!}
   \left|\varepsilon_{m-s}^{(2,\mathrm{mid})}\right|
   &\le\frac{234}{m^2},                                      \label{eq:audit-P3b}\\
 \sum_{s\le m/3}\frac{z_*^s}{s!}
   \left|\varepsilon_{m-s}^{(\ge3)}\right|
   &\le\frac{573}{m^2},                                      \label{eq:audit-P3c}\\
 \sum_{s\le m/3}\frac{z_*^s}{s!}\mathcal X_s
   &\le\frac{784}{m^2}.                                      \label{eq:audit-P4}
\end{align}
Here $\mathcal X_s$ is exactly the right-hand side of
\eqref{eq:cross-product-bound}, after the three pieces of $\varepsilon_{m-s}$
have been substituted.  The inequalities use only the following elementary
inputs:
\begin{itemize}[leftmargin=2em]
 \item $1-D_s\le(28/25)s/m^2$ on $3s\le m$;
 \item
 $\log\Pi_s\le(146/125)e_1(s)/m< (2237/10000)s$;
 \item the block estimate~\eqref{eq:Delta-term}, with the endpoint
 two-block compositions removed where appropriate;
 \item finite Poisson-moment bounds for
 $\sum z_*^s/s!$, $\sum sz_*^s/s!$, and the corresponding quadratic and
 quartic moments.
\end{itemize}
For~\eqref{eq:audit-P1}, the exponential remainder inequality
$0\le e^x-1-x\le x^2e^x/2$ is applied to $x=\log\Pi_s$; the tilt
$e^{(2237/10000)s}$ is absorbed by replacing $z_*$ with a fixed larger
rational parameter.  For~\eqref{eq:audit-P3a}, the exact endpoint ratio
$Nc_2c_{p-2}/c_p$ is compared with $5N/(36m^2)$ before the Poisson sum is
taken.  The two remaining block estimates follow from the same composition
bounds used in the residual envelope.  Finally,
\eqref{eq:audit-P4} is obtained by multiplying the already stated pointwise
bounds, rather than introducing a new analytic estimate.

The numerical Poisson bounds are finite rational certificates.  To verify one
of them, substitute $z_*=50/27$, truncate after the last degree that occurs,
clear the positive factorial denominators, and transform the resulting
polynomial in the auxiliary endpoint parameter to the Bernstein basis on
$[0,1]$.  Every coefficient is nonnegative.  This procedure also certifies the
two polynomial inequalities in Lemma~\ref{lem:near-base}; it is a finite
arithmetic check, not an appeal to numerical sampling.

It remains to account for the far signed tail.  Taking absolute values in the
last sum of~\eqref{eq:signlock-three}, applying the rational saddle bound to
$E^-_{m-s}(N)$, and summing the resulting Poisson tail from
$s=\floor{m/3}+1$ gives
\begin{equation}\label{eq:audit-far-signlock}
 |\mathcal R_{\mathrm{far}}|\le\frac1{m^2}.
\end{equation}
The first omitted term dominates the tail because its adjacent quotient is
less than one; after clearing factorials, the endpoint inequality is checked
at $m=361$ and is monotone thereafter.  Equations
\eqref{eq:audit-P1}--\eqref{eq:audit-far-signlock} sum to
\[
 \frac{426+13+184+234+573+784+1}{m^2}
 =\frac{2215}{m^2},
\]
which proves Proposition~\ref{prop:signlock-error-audit}.

\section{Scalar estimates in the Proposition~5.2 Gamma tail}
\label{app:gamma-tail-audit}

This appendix supplies the quantitative interfaces used in
Section~\ref{subsec:gamma52}.  Its organization mirrors the proof: first the
three coefficient errors in the exact split, then the Gamma moment entering
Jensen's inequality, and finally the four Taylor-remainder terms.

\subsection{The three split errors}

Put
\[
 \alpha_{a,s}=\frac{h_{a-s}}{Nc_a}
 \qquad(0\le s\le r_0).
\]
The factorial part of $c_{a-s}/c_a$ is the Gamma weight $\gamma_s$ from
\eqref{eq:gamma-weight52}; the remaining factor is
$D_s=d_{a-s}/d_a$, together with $1-s_{a-s}/N$.  The coefficientwise positive
majorant for $W$ therefore gives the exact aggregate estimate
\begin{equation}\label{eq:audit-H-error}
 \sum_{s=0}^{r_0}|\alpha_{a,s}-\gamma_s|\,|\omega_s|
 \le \frac{86069}{3125a^2}.
\end{equation}
The constant has a useful decomposition
\[
 \frac{86069}{3125}=36\left(\frac{2304}{3125}\right)+1.
\]
Here $2304/3125$ is the increment constant in the $d$-ratio estimate, the
factor~$36$ is the closed factorial-moment sum, and the final~$1$ covers the
partition-dependent factor $s_{a-s}/N$.  Thus
\eqref{eq:audit-H-error} is exactly the first inequality of
\eqref{eq:printed-errors52}, not a rounded asymptotic estimate.

For the other two terms, let $\widehat E$ and $\widehat W$ be obtained from
$E$ and $W$ by taking absolute values coefficientwise.  Positivity gives
\[
 |e_s|\le[t^s]\exp(\Lambda_+(t)),
 \qquad
 |\omega_s|\le[t^s]\exp(\Lambda_+(t))(1+J(t)),
\]
where $\Lambda_+(t)=\sum_{r\le p}h_rt^r$.  Evaluating these majorants at the
rational saddle radii used in Section~\ref{sec:saddle}, and then applying the
factorial-mass bounds there, yields
\begin{align}
 \frac1{Nc_a}\sum_{s=0}^{r_0}k_{a-s}|e_s|
 &\le\frac1{a^2},                                             \label{eq:audit-K-error}\\
 \frac{|\omega_a|}{Nc_a}
 &\le\frac1{a^2}.                                             \label{eq:audit-Omega-error}
\end{align}
The endpoint inequalities in both cases are polynomial after multiplication
by the positive factorial denominators.  Their shifted forms have
nonnegative coefficients in $a-150$, so the check at $a=150$ propagates to
all larger~$a$.

\subsection{The exponent moment in the Jensen step}

Let $G_{a-2}$ have Gamma shape $a-2$ and rate one and put
$T=1/(6G_{a-2})$.  Separating the $r=1$ term from the remaining terms in
$\Lambda$ gives
\begin{align}
 \mathbb E_{a-2}[\Lambda(T)]
 &\le \frac{5M}{24(a-3)} \\
 &\quad+
\frac{8M}{25}
 \left(
   \frac1{(a-3)(a-4)}
   +\frac{p-2}{3\binom{a-3}{3}}
 \right).                                                     \label{eq:audit-gamma-moment-raw}
\end{align}
The first line uses $h_1\le5M/4$ and the first negative Gamma moment.  For
$r\ge2$, the coefficient bound
$c_r\le(4/25)6^r(r-1)!$ cancels the factor $6^{-r}$ in the Gamma moment; the
remaining reciprocal falling factorials are bounded by their first term and a
binomial tail.  Since $p\le a/2$ and
\[
 \binom{a-3}{3}=\frac{(a-3)(a-4)(a-5)}6,
\]
the right side of~\eqref{eq:audit-gamma-moment-raw} is at most
\begin{equation}\label{eq:audit-gamma-moment-closed}
 \frac54+
\frac5{2(a-3)}+
\frac{96a^2}{25(a-5)^3}.
\end{equation}
After positive denominators are cleared, the assertion that
\eqref{eq:audit-gamma-moment-closed} is smaller than $13/10$ becomes
\[
 100a^4-14480a^3+110040a^2-410000a+662500>0.
\]
Writing this polynomial in powers of $a-150$ gives
\[
 100(a-150)^4+45520(a-150)^3+7094040(a-150)^2
 +405202000(a-150)+4170062500,
\]
so it is positive for every $a\ge150$.  This proves
\eqref{eq:gamma-exponent-bound} with no numerical approximation.

The other scalar endpoint in the Jensen step is
$e^{-13/10}>27/100$.  Equivalently, $e^{13/10}<100/27$; a finite positive
Taylor bound for $e^{3/10}$, multiplied by the standard rational upper bound
for $e$, proves this strict inequality.  Together with the Gamma shape shift
in~\eqref{eq:gamma-shape-shift}, this yields the margin
$9/(40(a-2))$.

\subsection{The four-term truncation residue}

Set $s_0=\floor{a/8}$ and let $R_a$ be the quantity in
\eqref{eq:truncation-residue52}.  Each term has a distinct origin:
\begin{center}
\begin{tabular}{p{0.43\textwidth}p{0.45\textwidth}}
\toprule
term & source\\
\midrule
$920/2^{r_0+1}$ & Taylor degrees above $r_0$ on $G\ge a/2$, where the
  evaluation radius is half the radius used for the coefficient majorant\\
$2(5/6)^a$ & the Gamma lower-tail event $G<a/2$, using
  $|e^{-\Lambda(T)}(1-J(T))|\le2$\\
$9(9/10)^{a-s_0}$ & retained Taylor degrees $s\le s_0$ on the lower event,
  controlled by the negative-Gamma-moment ratio\\
$920a(3/10)^{s_0+1}$ & degrees $s>s_0$, controlled by the same coefficient
  majorant at the larger radius\\
\bottomrule
\end{tabular}
\end{center}
After the coefficient bounds for $\Lambda$ and $J$ are inserted, the exact
endpoint certificate for the majorant is
\[
 \widehat W\!\left(\frac{2}{3a}\right)\le920.
\]  The
Gamma Chernoff estimate gives
$\mathbb P(G<a/2)\le(5/6)^a$.  The low-degree moment comparison contributes
the ratio $9/10$, while the radius comparison in the high degrees contributes
$3/10$.

For completeness, the inequality $R_a<a^{-2}$ is propagated in steps of
eight, which preserves $s_0=\floor{a/8}$ modulo its initial class.  After
multiplication by $a^2$, each of the four summands is nonincreasing under
$a\mapsto a+8$ for $a\ge150$.  Exact rational evaluation at
$a=150,151,\ldots,157$ gives a value below~$1$ in each residue class.  Hence
\[
 a^2R_a<1\qquad(a\ge150),
\]
which is the final inequality in~\eqref{eq:truncation-residue52}.

\section{Exact ratio and reserve estimates for the large positive tail}
\label{app:large-tail-details}

This appendix supplies the scalar inequalities behind the ratio--reserve
argument of the large positive tail.  Its purpose is twofold: to record the exact
inequalities and certificate interfaces used by the Lean proof of
Section~\ref{sec:positive}, so that the large positive tail can be audited
together with the formal certificate files, and to record exactly where the finite
prefix certificates enter.
All quantities below are rational.  We use the convention $0^0=1$.

\subsection{Candidate summands}

Fix $a>2000$.  Put
\begin{equation}\label{eq:app-tail-edges}
 N_-=6a-7,\qquad N_+=12a-8,\qquad
 s_- = \ceil{\sqrt{N_-}},\qquad s_+=\ceil{\sqrt{N_+}},
\end{equation}
and
\begin{equation}\label{eq:app-tail-indices}
 K=\floor{\frac{9a}{10}},\qquad
 \ell=\max\{1,s_-+1\},\qquad
 m=\floor{\frac a3}+10,\qquad j_k=a-k.
\end{equation}
For $a>2000$ one has
\begin{equation}\label{eq:app-tail-order}
 1\le \ell\le m<K<a.
\end{equation}
The entropy-shadow factor is
\begin{equation}\label{eq:app-tail-H}
 \mathcal H_a(k)=
 \frac{(k-1)^{k-1}(j_k-1)^{j_k-1}}{(a-2)^{a-2}}.
\end{equation}
The maximal-term estimate for the binomial expansion gives
\begin{equation}\label{eq:app-tail-binom-shadow}
 \frac{1}{(a-1)\binom{a-2}{k-1}}
 \le \mathcal H_a(k)
 \qquad (1\le k<a-1).
\end{equation}
For $k=1$ this is simply $1/(a-1)\le1$.  For $k\ge2$, put
$n=a-2$ and $r=k-1$.  In the binomial expansion
\[
 n^n=\sum_{i=0}^{n}\binom ni r^i(n-r)^{n-i},
\]
the summand at $i=r$ is maximal.  Hence
\[
 n^n\le(n+1)\binom nr r^r(n-r)^{n-r},
\]
which rearranges exactly to \eqref{eq:app-tail-binom-shadow}.

Recall the rational exponential upper bound
\begin{equation}\label{eq:app-tail-PE}
 \PE_T(x)=\sum_{q=0}^{T-1}\frac{x^q}{q!}
 +\frac{x^T}{T!}\frac{1}{1-x/T},
 \qquad 0\le x<T.
\end{equation}
Write
\begin{align}
 \Xi_s(a,k)
  &=\frac{1139}{1000}s_+
    +\frac{j_k}{5}+\frac{29}{10}\frac{a}{j_k}+1,
 \label{eq:app-tail-Es}\\
 \Xi_t(a,k)
  &=\frac a5+\frac{57}{10}\frac{a}{k}
    +\frac{29}{10}\frac{a}{j_k}+2,
 \label{eq:app-tail-Et}
\end{align}
and set
\begin{equation}\label{eq:app-tail-LsLt}
 L_s(a,k)=\PE_a(\Xi_s(a,k)),\qquad
 L_t(a,k)=\PE_{8a}(\Xi_t(a,k)).
\end{equation}
The two candidate summands used in the formal tail proof are
\begin{align}
 S_a(k)
  &=\frac{65}{N_+}\,k j_k\,\mathcal H_a(k)\,2^{-j_k}L_s(a,k),
 \label{eq:app-tail-S}\\
 T_a(k)
  &=\frac{96}{N_-}\,k j_k\,\mathcal H_a(k)\,2^{-j_k}L_t(a,k).
 \label{eq:app-tail-T}
\end{align}
They are not heuristic asymptotic expressions: they are the exact rational
majorants summed by the Lean proof.

To relate them to the normalized convolution in \eqref{eq:Unorm}, define
\begin{equation}\label{eq:app-tail-C}
 \mathcal C_a(N,k)=
 \frac N2 R_{k,a}2^{-j_k}\max\{X_k(N),0\}\,Y_{j_k}(N).
\end{equation}
The direct saddle estimates, the bound \eqref{eq:R-bound}, and
\eqref{eq:app-tail-binom-shadow} imply
\begin{align}
 \mathcal C_a(N,k)&\le S_a(k)
 &&\text{if }k\le\ceil{\sqrt N},
 \label{eq:app-tail-pointwise-small}\\
 \mathcal C_a(N,k)&\le T_a(k)
 &&\text{if }\ceil{\sqrt N}<k.
 \label{eq:app-tail-pointwise-temp}
\end{align}
The sign lock removes all $k>K$.  Consequently, for every
$N_-\le N\le N_+$,
\begin{equation}\label{eq:app-tail-product-to-two-sums}
 \sum_{1\le k<a}\mathcal C_a(N,k)
 \le \Sigma_s(a)+\Sigma_t(a),
\end{equation}
where
\begin{equation}\label{eq:app-tail-sums}
 \Sigma_s(a)=\sum_{k=1}^{\min\{K,s_+\}}S_a(k),
 \qquad
 \Sigma_t(a)=\sum_{k=\ell}^{K}T_a(k).
\end{equation}
The overlap is deliberately counted twice.  This loses a negligible amount
and allows the two regimes to be summed independently.

\subsection{The common adjacent quotient}

Let $1\le r<K$ and write $j=a-r$.  Removing the constants and exponential
shells from either \eqref{eq:app-tail-S} or \eqref{eq:app-tail-T} gives the
same adjacent quotient
\begin{equation}\label{eq:app-tail-qraw}
 \mathfrak q_a(r)=
 2\,
 \frac{(r+1)(j-1)r^r(j-2)^{j-2}}
 {rj(r-1)^{r-1}(j-1)^{j-1}}.
\end{equation}
Thus
\begin{equation}\label{eq:app-tail-ratio-exact}
 \frac{S_a(r+1)}{S_a(r)}
  =\mathfrak q_a(r)\frac{L_s(a,r+1)}{L_s(a,r)},
 \qquad
 \frac{T_a(r+1)}{T_a(r)}
  =\mathfrak q_a(r)\frac{L_t(a,r+1)}{L_t(a,r)}.
\end{equation}
For $r\ge2$ (and hence, in the retained range, $j\ge3$), it is useful to factor
\begin{equation}\label{eq:app-tail-P-factor}
 \mathfrak q_a(r)=2\frac{r+1}{j}\,\mathcal P_a(r),
\end{equation}
where
\begin{align}
 \mathcal P_a(r)
 &:=\frac{r^r}{r(r-1)^{r-1}}
 \frac{(j-1)(j-2)^{j-2}}{(j-1)^{j-1}}
 \notag\\
 &=\left(1+\frac1{r-1}\right)^{r-1}
   \left(1+\frac1{j-2}\right)^{-(j-2)}.
 \label{eq:app-tail-P}
\end{align}
Since $n\mapsto(1+1/n)^n$ is increasing, one obtains
\begin{equation}\label{eq:app-tail-P-directions}
 \mathcal P_a(r)\le1\quad\text{if }r-1\le j-2,
 \qquad
 \mathcal P_a(r)\ge1\quad\text{if }j-2\le r-1.
\end{equation}

We repeatedly use the elementary geometric estimates
\begin{align}
 F_{r+1}\le qF_r\ (u\le r<v)
 &\Longrightarrow
 \sum_{r=u}^{v}F_r\le\frac{F_u}{1-q},
 \label{eq:app-tail-forward-geom}\\
 F_{r-1}\le qF_r\ (u<r\le v)
 &\Longrightarrow
 \sum_{r=u}^{v}F_r\le\frac{F_v}{1-q},
 \label{eq:app-tail-reverse-geom}
\end{align}
valid for nonnegative $F_r$ and $0\le q<1$.

\subsection{The small branch}

\begin{lemma}[Small adjacent ratio]\label{lem:app-tail-small-ratio}
For $a>2000$ and
\(
 1\le r<\min\{K,s_+\},
\)
one has
\begin{equation}\label{eq:app-tail-small-ratio}
 S_a(r+1)\le\frac12S_a(r).
\end{equation}
\end{lemma}

\begin{proof}
Since $s_+=\ceil{\sqrt{12a-8}}$, an elementary integer estimate gives
\begin{equation}\label{eq:app-tail-small-cutoff}
 12s_+\le a\qquad(a>2000),
\end{equation}
because $\ceil{\sqrt{12a-8}}<\sqrt{12a-8}+1$ and
\[
 (a-12)^2-144(12a-8)=a^2-1752a+1296>0\qquad(a\ge2001)
\]
(positive at $a=2001$ and increasing for $a>876$), so
$12\bigl(\sqrt{12a-8}+1\bigr)\le a$.  As $r<s_+$, hence
$r+1\le s_+$, this gives $12(r+1)\le a$; writing $j=a-r$ then yields
$j\ge11(r+1)+1$, and therefore
\begin{equation}\label{eq:app-tail-small-linear-gap}
 272(r+1)\le25j.
\end{equation}
For $r\ge2$ the factorization
\eqref{eq:app-tail-P-factor}--\eqref{eq:app-tail-P} bounds the first power ratio
by $\e<\tfrac{68}{25}$ (by the classical bound $(1+1/n)^n<\e$) and the reciprocal
second ratio by $1$, so $\mathcal P_a(r)\le\tfrac{68}{25}$ and
\begin{equation}\label{eq:app-tail-qraw-half}
 \mathfrak q_a(r)=2\frac{r+1}{j}\,\mathcal P_a(r)
 \le\frac{136(r+1)}{25j}\le\frac12
\end{equation}
by \eqref{eq:app-tail-small-linear-gap}.  For $r=1$ one has directly
\[
 \mathfrak q_a(1)=\frac4j\Bigl(\frac{j-2}{j-1}\Bigr)^{j-2}\le\frac4j\le\frac12,
\]
since $j=a-1>1999$, so \eqref{eq:app-tail-qraw-half} holds throughout
$1\le r<\min\{K,s_+\}$.  The small exponent is nonincreasing: with $j=a-r$,
\eqref{eq:app-tail-Es} gives
\[
 \Xi_s(a,r+1)-\Xi_s(a,r)=-\frac15+\frac{29a}{10\,j(j-1)}\le0,
\]
since $12(r+1)\le a$ yields $j\ge\tfrac{11a}{12}+1$ and hence $29a\le2j(j-1)$, so
\begin{equation}\label{eq:app-tail-Es-mono}
 \Xi_s(a,r+1)\le \Xi_s(a,r),
\end{equation}
and monotonicity of \eqref{eq:app-tail-PE} gives $L_s(a,r+1)\le L_s(a,r)$.
Substituting \eqref{eq:app-tail-qraw-half} and this bound into
\eqref{eq:app-tail-ratio-exact} proves \eqref{eq:app-tail-small-ratio}.
\end{proof}

The first candidate has the exact form
\begin{equation}\label{eq:app-tail-S1}
 S_a(1)=\frac{65(a-1)}{12a-8}\,2^{-(a-1)}L_s(a,1),
\end{equation}
because $\mathcal H_a(1)=1$.  The exponential shell satisfies
\begin{equation}\label{eq:app-tail-small-shell}
 L_s(a,1)\le\left(\frac32\right)^a.
\end{equation}
Indeed $\Xi_s(a,1)\le3a/10$, and the finite rational exponential bound at
$3a/10$ is dominated by the negative-binomial envelope $(3/2)^a$.
Therefore
\begin{equation}\label{eq:app-tail-small-envelope}
 800000000\,S_a(1)
 \le
 800000000\frac{65(a-1)}{12a-8}
 2^{-(a-1)}\left(\frac32\right)^a.
\end{equation}
The right-hand side is nonincreasing for $a\ge3$; its value at $a=2001$ is
at most one by an exact rational computation.  Hence
\begin{equation}\label{eq:app-tail-small-first-unit}
 800000000\,S_a(1)\le1\qquad(a>2000).
\end{equation}
Combining Lemma~\ref{lem:app-tail-small-ratio},
\eqref{eq:app-tail-forward-geom}, and
\eqref{eq:app-tail-small-first-unit} gives
\begin{equation}\label{eq:app-tail-small-budget}
 \Sigma_s(a)\le2S_a(1)\le\frac1{400000000}.
\end{equation}
The actual guarded convolution term at $k=1$ vanishes because $X_1=-1$.
The nonzero value $S_a(1)$ is retained only as a convenient first-term reserve
for the candidate sequence.

\subsection{The lower tempered band}

Put
\begin{equation}\label{eq:app-tail-q}
 q_a=\frac{4a-1}{4a}=1-\frac1{4a}.
\end{equation}
The lower tempered band is $\ell\le k\le m$.  Its desired step estimate is
\begin{equation}\label{eq:app-tail-temp-lower-step}
 T_a(r+1)\le q_aT_a(r)
 \qquad(\ell\le r<m).
\end{equation}
On this range $r-1\le j-2$, so
$\mathcal P_a(r)\le1$ by \eqref{eq:app-tail-P-directions}, and therefore
\begin{equation}\label{eq:app-tail-qraw-lower}
 \mathfrak q_a(r)\le2\frac{r+1}{j}.
\end{equation}
Define the sharp exponential quotient target
\begin{equation}\label{eq:app-tail-U}
 U^{\sharp}_a(r)=\frac12\frac{j}{r+1}q_a.
\end{equation}
Then \eqref{eq:app-tail-qraw-lower} gives the exact budget identity
\begin{equation}\label{eq:app-tail-U-budget}
 \mathfrak q_a(r)U^{\sharp}_a(r)\le q_a.
\end{equation}
Thus it remains only to prove
\begin{equation}\label{eq:app-tail-Lt-forward-target}
 \frac{L_t(a,r+1)}{L_t(a,r)}\le U^{\sharp}_a(r).
\end{equation}

There are two subranges.  If
\begin{equation}\label{eq:app-tail-front-range}
 3(r+1)\le a,
\end{equation}
then
\begin{equation}\label{eq:app-tail-Et-front-mono}
 \Xi_t(a,r+1)\le \Xi_t(a,r),
 \qquad U^{\sharp}_a(r)\ge1.
\end{equation}
Hence \eqref{eq:app-tail-Lt-forward-target} follows immediately from
monotonicity of $\PE_{8a}$.

The only remaining indices are the ten offsets
\begin{equation}\label{eq:app-tail-top-offset}
 r=\floor{\frac a3}+t,
 \qquad 0\le t<10.
\end{equation}
For $a\ge3000$ the following two rational inequalities hold uniformly in
$t$:
\begin{align}
 \Xi_t(a,r+1)+\frac{175}{4a}&\le \Xi_t(a,r),
 \label{eq:app-tail-top-drop}\\
 1&\le U^{\sharp}_a(r)\left(1+\frac{175}{12a}\right)^3.
 \label{eq:app-tail-top-budget}
\end{align}
The second estimate is reduced to the nonnegativity of
\begin{equation}\label{eq:app-tail-top-poly}
 550656a^3-76358800a^2-559763750a+144703125.
\end{equation}
After writing $a=3000+x$, this polynomial becomes
\begin{equation}\label{eq:app-tail-top-poly-shift}
 550656x^3+4879545200x^2
 +14408999436250x+14178803653453125,
\end{equation}
which is manifestly nonnegative for $x\ge0$.

The finite-exponential three-step shift inequality says that, whenever
$y+d\le z$,
\begin{equation}\label{eq:app-tail-three-step}
 \PE_T(y)\le
 \left(1+\frac d3\right)^{-3}\PE_T(z)
\end{equation}
provided the displayed arguments lie in the certified range.  Applying it
with $d=175/(4a)$ and using
\eqref{eq:app-tail-top-drop}--\eqref{eq:app-tail-top-budget} proves
\eqref{eq:app-tail-Lt-forward-target} for every $a\ge3000$.

For the finite prefix $2001\le a<3000$, define
\begin{equation}\label{eq:app-tail-theta}
 \theta_a=1-\frac{213}{5a}=\frac{5a-213}{5a}.
\end{equation}
Exact rational certificates verify, for every guarded pair
$(a,t)$ with \eqref{eq:app-tail-top-offset},
\begin{align}
 \frac{L_t(a,r+1)}{L_t(a,r)}&\le\theta_a,
 \label{eq:app-tail-prefix-exp-ratio}\\
 (4a)\,\mathfrak q_a(r)\theta_a&\le4a-1.
 \label{eq:app-tail-prefix-raw}
\end{align}
The first computation is split into twenty blocks of fifty consecutive values
of $a$, each containing the ten offsets $t=0,\ldots,9$.  The second is a
single $999\times10$ raw-rational block.  These are exact Boolean checks over
$\mathbb Q$ executed by Lean's native evaluator; no floating-point comparison
is used.  Equations \eqref{eq:app-tail-prefix-exp-ratio} and
\eqref{eq:app-tail-prefix-raw} imply \eqref{eq:app-tail-temp-lower-step} in the
prefix strip.  Together with the symbolic argument above, they prove
\eqref{eq:app-tail-temp-lower-step} for every $a>2000$.

\subsection{The upper tempered band}

The upper band is $m+1\le k\le K$ and is summed in reverse.  The desired step
is
\begin{equation}\label{eq:app-tail-temp-upper-step}
 T_a(r-1)\le q_aT_a(r)
 \qquad(m+1<r\le K).
\end{equation}
Set
\begin{equation}\label{eq:app-tail-V}
 V^{\sharp}_a(r)=\mathfrak q_a(r-1)q_a.
\end{equation}
By the exact quotient identity \eqref{eq:app-tail-ratio-exact}, it is enough
to prove
\begin{equation}\label{eq:app-tail-Lt-reverse-target}
 L_t(a,r-1)\le V^{\sharp}_a(r)L_t(a,r).
\end{equation}

In the far upper band
\begin{equation}\label{eq:app-tail-far-upper}
 3a\le5r,
\end{equation}
one has
\begin{equation}\label{eq:app-tail-far-upper-two}
 \Xi_t(a,r-1)\le \Xi_t(a,r),
 \qquad V^{\sharp}_a(r)\ge1.
\end{equation}
The second inequality follows from the lower form of
\eqref{eq:app-tail-P-directions}, applied to the raw quotient at $r-1$.
Thus \eqref{eq:app-tail-Lt-reverse-target} is immediate in this range.

It remains to treat the middle band
\begin{equation}\label{eq:app-tail-middle-upper}
 m+1<r\le K,
 \qquad 5r<3a.
\end{equation}
The exponent shift is bounded by
\begin{equation}\label{eq:app-tail-middle-shift}
 \Xi_t(a,r-1)\le \Xi_t(a,r)+\delta_a,
 \qquad \delta_a=\frac{45}{a}.
\end{equation}
Put
\begin{equation}\label{eq:app-tail-F}
 \Theta_a=\left(1-\frac{79}{75}\delta_a\right)^{-1}.
\end{equation}
The rational exponential estimate used in Lean is
\begin{equation}\label{eq:app-tail-shift-difference}
 \PE_{8a}(E+\delta_a)-\PE_{8a}(E)
 \le \frac{79}{75}\delta_a\,\PE_{8a}(E+\delta_a),
\end{equation}
valid here because
\begin{equation}\label{eq:app-tail-shift-cutoff}
 E+\delta_a\le\frac{8a-1}{24}.
\end{equation}
Rearranging \eqref{eq:app-tail-shift-difference} and using
\eqref{eq:app-tail-middle-shift} gives
\begin{equation}\label{eq:app-tail-middle-shell}
 L_t(a,r-1)\le \Theta_aL_t(a,r).
\end{equation}
It remains to show $\Theta_a\le V^{\sharp}_a(r)$.

For this purpose write the raw quotient at $r-1$ as
\begin{equation}\label{eq:app-tail-reverse-factor}
 \mathfrak q_a(r-1)=
 \frac{2r}{a-r+1}\,\mathcal P_a(r-1).
\end{equation}
Define
\begin{equation}\label{eq:app-tail-W}
 W_a(r)=\Theta_a\frac{4a}{4a-1}\frac{a-r+1}{2r}.
\end{equation}
Then $\Theta_a\le V^{\sharp}_a(r)$ is equivalent to
\begin{equation}\label{eq:app-tail-W-target}
 W_a(r)\le\mathcal P_a(r-1).
\end{equation}
There are two elementary subcases.

If $a+1\le2r$, then direct rational simplification gives
\begin{equation}\label{eq:app-tail-W-midpoint}
 W_a(r)\le1,
\end{equation}
and \eqref{eq:app-tail-P-directions} gives
$1\le\mathcal P_a(r-1)$.

If $2r<a+1$, then
\begin{equation}\label{eq:app-tail-W-lowhalf}
 W_a(r)\le1-\frac1a.
\end{equation}
The required matching lower bound for the power product is obtained from the
telescoping inequality
\begin{equation}\label{eq:app-tail-P-gap}
 \mathcal P_a(s)\ge
 1-\frac{(j_s-2)-(s-1)}{2(s-1)(j_s-2)}
 \qquad(s-1\le j_s-2).
\end{equation}
For $s=r-1$ and the present range, the fraction on the right is at most
$1/a$; hence
\begin{equation}\label{eq:app-tail-P-lowhalf}
 1-\frac1a\le\mathcal P_a(r-1).
\end{equation}
This proves \eqref{eq:app-tail-W-target}, then
\eqref{eq:app-tail-Lt-reverse-target}, and finally
\eqref{eq:app-tail-temp-upper-step} throughout the upper band.

\subsection{Endpoint reserves and the tempered budget}

The two endpoint exponential shells admit the common envelope
\begin{equation}\label{eq:app-tail-ten-sevenths}
 L_t(a,\ell)\le\left(\frac{10}{7}\right)^a,
 \qquad
 L_t(a,K)\le\left(\frac{10}{7}\right)^a.
\end{equation}
The theorem-facing endpoint estimates are
\begin{equation}\label{eq:app-tail-endpoint-units}
 800000000\,(4a)T_a(\ell)\le1,
 \qquad
 800000000\,(4a)T_a(K)\le1.
\end{equation}
We record the coarse envelopes used to prove them.

For the lower endpoint, one has $\ell\le a/4$, discards
$\mathcal H_a(\ell)\le1$, and keeps only the dyadic decay.  This yields
\begin{equation}\label{eq:app-tail-Gminus}
 800000000(4a)T_a(\ell)
 \le G_-(a):=
 3200000000\,a^3\,2^{-(a-\floor{a/4})}
 \left(\frac{10}{7}\right)^a.
\end{equation}
The four-step comparison
\begin{equation}\label{eq:app-tail-Gminus-step}
 G_-(a+4)\le G_-(a)
\end{equation}
follows from
\begin{equation}\label{eq:app-tail-Gminus-constants}
 a+4\le\frac{401}{400}a,
 \qquad
 10000\left(\frac{401}{400}\right)^3\le19208.
\end{equation}
Exact rational checks at $a=2001,2002,2003,2004$ initialize the four residue
classes and prove the first inequality in
\eqref{eq:app-tail-endpoint-units}.

For the upper endpoint, the entropy shadow is retained.  The elementary
bounds
\begin{equation}\label{eq:app-tail-upper-entropy}
 \mathcal H_a(K)
 \le
 \left(\frac{901}{1000}\right)^{\floor{89a/100}}
 \left(\frac{101}{1000}\right)^{\floor{a/10}-1}
\end{equation}
give
\begin{align}
 800000000(4a)T_a(K)
 &\le G_+(a)\notag\\
 &:={}
 3200000000\,a^3
 \left(\frac{901}{1000}\right)^{\floor{89a/100}}
 \left(\frac{101}{1000}\right)^{\floor{a/10}-1}
 2^{-\floor{a/10}}
 \left(\frac{10}{7}\right)^a.
 \label{eq:app-tail-Gplus}
\end{align}
The one-hundred-step comparison
\begin{equation}\label{eq:app-tail-Gplus-step}
 G_+(a+100)\le G_+(a)
\end{equation}
follows from
\begin{equation}\label{eq:app-tail-Gplus-constants}
 \left(\frac{21}{20}\right)^3
 \left(\frac{901}{1000}\right)^{89}
 \left(\frac{101}{1000}\right)^{10}
 \left(\frac{10}{7}\right)^{100}
 \le2^{10}.
\end{equation}
Exact rational checks for $2001\le a\le2100$ initialize the one hundred
residue classes and prove the second inequality in
\eqref{eq:app-tail-endpoint-units}.

Since $1-q_a=1/(4a)$, the lower and upper geometric reserves satisfy
\begin{align}
 \sum_{k=\ell}^{m}T_a(k)
 &\le\frac{T_a(\ell)}{1-q_a}
  =4aT_a(\ell)\le\frac1{800000000},
 \label{eq:app-tail-lower-budget}\\
 \sum_{k=m+1}^{K}T_a(k)
 &\le\frac{T_a(K)}{1-q_a}
  =4aT_a(K)\le\frac1{800000000}.
 \label{eq:app-tail-upper-budget}
\end{align}
Therefore
\begin{equation}\label{eq:app-tail-tempered-budget}
 \Sigma_t(a)\le\frac1{400000000}.
\end{equation}
Together with \eqref{eq:app-tail-small-budget}, this gives the product part of
the positive budget:
\begin{equation}\label{eq:app-tail-product-budget}
 \Sigma_s(a)+\Sigma_t(a)\le\frac1{200000000}
 \qquad(a>2000).
\end{equation}

\subsection{The solo term}

The remaining term in \eqref{eq:Unorm} is $2^{-a-1}Y_a(N)$.  The sharp solo
estimate is proved first at the upper edge $N=N_+$ by splitting the closed
factorial-composition expansion into low-degree and remainder blocks.  In the
notation of the formal development, the closed block sum is
\begin{equation}\label{eq:app-tail-solo-blocksum}
\begin{aligned}
 \Delta_a(N)
 ={}&\sum_{s=0}^{a-1}\frac{(Nc_1/4)^s}{s!}
 \sum_{\substack{r\ge1\\2r\le a-s}}
 \frac{(2N/25)^r\,3^{\,a-s}\,4^{\,r-1}\,(a-s-2r+1)!}{r!}\\
 &+\frac{(Nc_1/4)^a}{a!},
\end{aligned}
\end{equation}
which is \code{positiveLargeTailSoloSharpGcompClosedFactorialSplitBlockSum} in
Lean.  The denominator-cleared estimate at the upper edge is
\begin{equation}\label{eq:app-tail-solo-sharp}
 4\,2^a\Delta_a(N_+)
 \le29\,a\,c_a\left(\frac{10}{7}\right)^a.
\end{equation}
Monotonicity in $N$ then yields the normalized bound
\begin{equation}\label{eq:app-tail-solo-envelope}
 Y_a(N)\le\frac{29}{2}\frac aN
 \left(\frac{10}{7}\right)^a
 \qquad(N_-\le N\le N_+).
\end{equation}
Since $N\ge6a-7$ and $a\ge401$, one has $a/N\le1/5$.  Therefore
\begin{align}
 2^{-a-1}Y_a(N)
 &\le\frac{29}{4}\frac aN\left(\frac57\right)^a
 \notag\\
 &\le\frac{29}{20}\left(\frac57\right)^a.
 \label{eq:app-tail-solo-geometric}
\end{align}
Writing
\[
 H_{\mathrm{solo}}(a)=290000000\left(\frac57\right)^a,
\]
one has the exact adjacent-$a$ relation
\begin{equation}\label{eq:app-tail-solo-step}
 H_{\mathrm{solo}}(a+1)=\frac57H_{\mathrm{solo}}(a).
\end{equation}
Thus the sequence is decreasing, and the exact rational check at $a=401$
gives
\begin{equation}\label{eq:app-tail-solo-unit}
 290000000\left(\frac57\right)^a\le1.
\end{equation}
Consequently
\begin{equation}\label{eq:app-tail-solo-budget-final}
 2^{-a-1}Y_a(N)\le\frac1{200000000}.
\end{equation}

\begin{proposition}[Large positive-tail budget]
\label{prop:app-tail-final}
For every $a>2000$ and every $N$ with $6a-7\le N\le12a-8$,
\begin{equation}\label{eq:app-tail-final}
 2^{-a-1}Y_a(N)
 +\frac N2\sum_{\substack{1\le k<a\\X_k(N)>0}}
 R_{k,a}2^{-(a-k)}X_k(N)Y_{a-k}(N)
 \le10^{-8}.
\end{equation}
\end{proposition}

\begin{proof}
The sign lock removes $k>K$.  The remaining product sum is at most
\eqref{eq:app-tail-product-budget} by
\eqref{eq:app-tail-pointwise-small}--\eqref{eq:app-tail-product-to-two-sums}.
The solo term is at most \eqref{eq:app-tail-solo-budget-final}.  Their sum is
\(
 1/200000000+1/200000000=10^{-8}.
\)
\end{proof}

\subsection{Correspondence with the formal proof}

For audit purposes, the principal Lean declarations corresponding to this
appendix are listed below.  The declaration names are stable in the release
version of the repository.

\begin{center}
\scriptsize
\urlstyle{tt}
\begin{tabular}{p{0.35\textwidth}p{0.57\textwidth}}
\toprule
paper item & Lean declaration or file\\
\midrule
$L_s,L_t$ &
\path{positiveSmallLargeExp}; \path{positiveTemperedLargeExp}\\
$S_a,T_a$ &
\path{positiveSmallEntropyShadowExpMajorantTerm};
\path{positiveTemperedEntropyShadowExpMajorantTerm}\\
raw quotient $\mathfrak q_a$ &
\path{positiveEntropyShadowBaseStepRawQuotient}\\
power product $\mathcal P_a$ &
\path{positiveEntropyShadowBaseStepRawPowerProduct}\\
small ratio $1/2$ &
\path{positiveSaddleLargeTailCandidateSmallRawBaseHalfCertificate}\\
lower split $m=\floor{a/3}+10$ &
\path{positiveLargeExpTemperedSplit}\\
lower sharp target $U^{\sharp}_a$ &
\path{positiveTemperedLowerSharpExpQuotientTarget}\\
finite prefix target $\theta_a$ &
\path{positiveTemperedLowerPrefixTopOffsetExpRatioTarget}\\
finite prefix certificates &
\path{PositiveSaddlePrefixExpRatioCertificate.lean};
\path{PositiveSaddlePrefixRawBudget.lean}\\
upper-middle shift $45/a$ &
\path{positiveTemperedUpperMiddleExpShiftBudget}\\
upper-middle factor $\Theta_a$ &
\path{positiveTemperedUpperMiddleExpShiftFactor}\\
small first reserve &
\path{positiveSaddleLargeTailCandidateSmallFirstReserveCertificate_closed}\\
tempered endpoint reserves &
\path{positiveSaddleLargeTailCandidateUnitReserveCertificate_temperedTenSevenths_closed}\\
assembled hybrid ratio certificate &
\path{positiveSaddleLargeTailCandidateRefinedAtomicCertificate_hybridClosed}\\
solo block sum $\Delta_a$ &
\path{positiveLargeTailSoloSharpGcompClosedFactorialSplitBlockSum}\\
solo block-sum cleared &
\path{positiveLargeTailSoloSharpGcompClosedFactorialSplitBlockSumTenSeventhsCleared}\\
solo saddle cleared &
\path{positiveLargeTailSoloSharpGcompSaddleTenSeventhsCleared}\\
solo scalar budget &
\path{positiveLargeTailSoloTenSeventhsScalarBudget_of_ge_401}\\
\bottomrule
\end{tabular}
\end{center}

\clearpage
\section{Public Lean statement}\label{app:lean-statement}

The following is a Lean-style rendering of the public facade, lightly edited for
readability (e.g.\ \code{Q[[X]]} for the rational power-series type and ASCII
\code{<=}, \code{!=} for the order and disequality); it is not literal source.
The exact, machine-checked file is \code{Prop51/Theorem.lean} in the repository.
It is reproduced to make the correspondence between the paper and the formal
theorem immediate.

\begin{lstlisting}
import Prop51.BCoeffSeries
import Prop51.Completion

namespace Prop51

open PowerSeries

noncomputable abbrev chenLarsonC : Q[[X]] := Cseries

theorem coeff_chenLarsonC (k : Nat) :
    coeff k chenLarsonC =
      (Nat.factorial (6*k) : Q) /
        ((Nat.factorial (3*k) : Q) *
          (Nat.factorial (2*k) : Q) * 72^k)

noncomputable abbrev chenLarsonSeries (mu : List Nat) : Q[[X]] :=
  bSeries mu

theorem chenLarsonSeries_spec (mu : List Nat) :
    chenLarsonC ^ ((mu.map (fun m => m + 1)).sum) *
      chenLarsonSeries mu
      = (mu.map (fun m =>
          rescale (((m + 1 : Nat) : Q)^(-1)) chenLarsonC)).prod

theorem chenLarsonCoefficient_neg
    {g : Nat} (hg : 2 <= g) (hmod : g % 3 != 1)
    {mu : List Nat} (hmu : IsPartitionOf mu (2 * g - 2)) :
    coeff (g / 3 + 1) (chenLarsonSeries mu) < 0

theorem chenLarsonCoefficient_ne_zero
    {g : Nat} (hg : 2 <= g) (hmod : g % 3 != 1)
    {mu : List Nat} (hmu : IsPartitionOf mu (2 * g - 2)) :
    coeff (g / 3 + 1) (chenLarsonSeries mu) != 0

end Prop51
\end{lstlisting}

\clearpage
For the corrected Proposition~5.2 the analogous public facade is
\code{Prop52/Theorem.lean}.  Besides re-exporting the series of
Proposition~5.1, it \emph{defines} the corrected source factor
$D_\mu^{\mathrm{cor}}$ coefficientwise, pins its coefficients, and states the two
final source theorems (again lightly edited for readability).

\begin{lstlisting}
import Prop51.Theorem
import Prop52.Source

namespace Prop52

-- B_mu(X) is re-exported from Prop 5.1.
noncomputable abbrev chenLarsonSeries (mu : List Nat) : Q[[X]] :=
  Prop51.chenLarsonSeries mu

-- The corrected factor D^cor, defined coefficientwise.
-- Geometrically lead = M(g/3+1) = 2g-2; printed factor uses 1.
noncomputable abbrev chenLarsonProp52SourceFactor
    (lead : Q) (mu : List Nat) : Q[[X]] :=
  mk (sourceFactorCoeff lead mu)

-- Its coefficients pin D^cor exactly:
theorem coeff_sourceFactor_zero (lead : Q) (mu : List Nat) :
    coeff 0 (chenLarsonProp52SourceFactor lead mu) = lead

theorem coeff_sourceFactor_one (lead : Q) (mu : List Nat) :
    coeff 1 (chenLarsonProp52SourceFactor lead mu)
      = -2 * ((N mu : Q) - sPower mu 1)

theorem coeff_sourceFactor_succ_succ
    (lead : Q) (mu : List Nat) (r : Nat) :
    coeff (r + 2) (chenLarsonProp52SourceFactor lead mu)
      = -12 * (r + 1) * Prop51.c (r + 1)
          * ((N mu : Q) - sPower mu (r + 2))

-- Non-vanishing for g >= 2, g = 1 (mod 3):
theorem chenLarsonProp52Coefficient_nonvanishing
    {g : Nat} (hg : 2 <= g) (hmod : g % 3 = 1)
    {mu : List Nat} (hmu : Prop51.IsPartitionOf mu (2 * g - 2)) :
    coeff (g / 3 + 1)
      (chenLarsonSeries mu
        * chenLarsonProp52SourceFactor (M (g / 3 + 1)) mu) != 0

-- ... and strict negativity for g >= 40:
theorem chenLarsonProp52Coefficient_neg
    {g : Nat} (hg : 40 <= g) (hmod : g % 3 = 1)
    {mu : List Nat} (hmu : Prop51.IsPartitionOf mu (2 * g - 2)) :
    coeff (g / 3 + 1)
      (chenLarsonSeries mu
        * chenLarsonProp52SourceFactor (M (g / 3 + 1)) mu) < 0

end Prop52
\end{lstlisting}

Here the non-vanishing facade carries \code{2 <= g} where
Theorem~\ref{thm:main52} states $g\ge4$; the two select the same genera, since
for $g\equiv1\pmod3$ the condition $g\ge2$ already forces $g\ge4$.  The
strict-negativity facade carries \code{40 <= g}, matching $g\ge40$ directly.

\section{Paper-to-Lean correspondence}\label{app:map}

\begin{center}
\small
\begin{tabular}{p{0.34\textwidth}p{0.56\textwidth}}
\toprule
paper item & principal Lean declaration or file\\
\midrule
hypergeometric series $\C$ & \code{Aseq}, \code{Cseries} in
  \code{ExpSeries.lean}\\
$\C=\exp(\sum c_rX^r)$ & \code{Cseries\_eq\_expSeries\_c} in
  \code{Bridge.lean}\\
quotient-series identity & \code{bSeries\_official} in
  \code{BCoeffSeries.lean}\\
partition-free majorant & \code{bCoeff\_le\_U} in
  \code{Majorant.lean}\\
coefficient bounds & \code{d\_lb}, \code{d\_ub},
  \code{d\_increment}, \code{d\_ratio\_lb} in \code{DNorm.lean}\\
sign lock & \code{Xnorm\_le\_neg\_final\_margin} in
  \code{SignLock.lean}\\
finite saddle and prefix convolution & \code{expCoeff\_saddle},
  \code{expPrefix\_mul\_le} in \code{SaddleDirect.lean}\\
direct product estimates & \code{Bq\_Qq\_small\_core},
  \code{Bq\_Qq\_tempered\_core} in \code{SaddleDirect.lean}\\
finite edge budget & \code{positiveSaddleFiniteScaledEdgeBudget} in
  \path{Prop51/Generated/PositiveSaddleFiniteScaledEdge.lean}\\
closed proof & \code{completion\_of\_directSaddle\_closed} and
  \code{coefficientNegativity}\\
public theorem & \code{chenLarsonCoefficient\_neg},
  \code{chenLarsonCoefficient\_ne\_zero} in \code{Prop51/Theorem.lean}\\
\midrule
\multicolumn{2}{l}{\emph{Corrected Proposition~5.2}}\\
corrected source coefficient $T_a^{\mathrm{cor}}$ & \code{sourceCoeff} in
  \code{Prop52/Source.lean}\\
source--marked bridge & \code{sourceCorrectedCoeff\_eq} in
  \code{Prop52/Source.lean}\\
central identity~\eqref{eq:central52} &
  \code{correctedCoeff\_eq\_printedCoeff\_add} in \code{Prop52/Assembly.lean}\\
rectangle quotient sign & \code{Prop51.bCoeff\_neg\_of\_rectangle} in
  \code{Prop51/Rectangle.lean}\\
finite range $2\le a\le13$ & \code{correctedCoeff\_finite\_nonvanishing} in
  \code{Prop52/Finite.lean}\\
printed sign, $14\le a\le149$ & \code{printedCoeffNegativityMid\_closed} in
  \code{Prop52/MidBridge.lean}\\
printed sign, Gamma tail & \code{printedTailERealSeriesHasSum} in
  \code{Prop52/GammaCompletion.lean}\\
public theorems (Prop.~5.2) &
  \code{chenLarsonProp52Coefficient\_nonvanishing},
  \code{chenLarsonProp52Coefficient\_neg} in \code{Prop52/Theorem.lean}\\
\bottomrule
\end{tabular}
\end{center}

\section{Computational certificate inventory}\label{app:certs}

\begin{center}
\small
\begin{tabular}{p{0.18\textwidth}p{0.23\textwidth}p{0.46\textwidth}}
\toprule
range & method & verification surface\\
\midrule
$a\le8$ & exact partition enumeration & exact rational coefficient for every
positive partition; selected partition-cardinality checks\\
$9\le a\le60$ & exact rational majorant & all $10764$ rectangle pairs\\
$61\le a\le400$ & $192$-bit outward dyadic intervals & proved enclosure of
recurrence tables; eight native chunks\\
$401\le a\le2000$ & direct saddle plus fixed-point edge certificate & sixteen
$100$-row shards, scale $10^{20}$; analytic solo theorem\\
$2001\le a<3000$ & direct saddle plus generated prefix ratios & pointwise
product and solo bounds; finite quotient/raw-budget chunks\\
$a\ge3000$ & symbolic rational tail & adjacent quotients, reverse upper band,
unit reserves, and sharp solo envelope\\
\bottomrule
\end{tabular}
\end{center}

\noindent For the corrected Proposition~5.2 the ranges are as follows.

\begin{center}
\small
\begin{tabular}{p{0.16\textwidth}p{0.30\textwidth}p{0.42\textwidth}}
\toprule
range & method & verification surface\\
\midrule
$2\le a\le8$ & exact rational partition enumeration & exact corrected coefficient
for every positive partition (with permutation/generator completeness)\\
$9\le a\le13$ & modular certificate at $p=1{,}000{,}000{,}007$ & nonzero residue,
plus the rational-to-$\mathbb F_p$ reduction bridge\\
$14\le a\le149$ & exact dyadic interval certificate & strict negativity of the
printed marked coefficient $T_a^{\mathrm{old}}$\\
$a\ge150$ & Gamma-integral, moment, and Taylor-truncation estimates & real
exponential power-series identity for the upper-event integrand\\
$a\ge14$ (assembly) & central identity~\eqref{eq:central52} with the rectangle &
$T_a^{\mathrm{cor}}<0$ from $T_a^{\mathrm{old}}<0$ and $b_a<0$\\
\bottomrule
\end{tabular}
\end{center}

\begin{thebibliography}{99}

\bibitem{ChenLarson}
D.~Chen and H.~Larson,
\emph{Independence of tautological classes and cohomological stability for
strata of differentials}, arXiv:2603.23850v1, 25~March~2026
(Proposition~5.2 and equation~(5.4)).

\bibitem{Ionel}
E.-N.~Ionel,
\emph{Relations in the tautological ring of $\mathcal M_g$},
Duke Math. J. \textbf{129} (2005), no.~1, 157--186; arXiv:math/0312100
(Proposition~1.7, equation~(1.13)).

\bibitem{LeanAxioms}
The Lean development team,
\emph{Lean Reference Manual: Axioms and validating proofs},
\url{https://lean-lang.org/doc/reference/latest/Axioms/}, accessed June 2026.

\bibitem{Mathlib}
The Mathlib community,
\emph{Mathlib: the mathematical library for Lean 4},
\url{https://github.com/leanprover-community/mathlib4}.

\bibitem{FengAletheia}
T.~Feng et al.,
\emph{Towards Autonomous Mathematics Research},
arXiv:2602.10177, 2026.

\bibitem{OpenAIGoals}
R.~Pathak and S.~Fabbri,
\emph{Using Goals in Codex}, OpenAI Developers, 9 May 2026,
\url{https://developers.openai.com/cookbook/examples/codex/using_goals_in_codex},
accessed June 2026.

\end{thebibliography}

\end{document}
